\section{Protocol details}\label{app:protocol}

\begin{table}[htbp]
\centering\footnotesize\renewcommand{\arraystretch}{0.9}
\caption{Real-engine stages (SLURM job IDs on Nova). Submit, start, end and node of every job: \fpath{provenance/sacct_cited_jobs_2026-11-01.txt} (local time; \fpath{_utc.txt}, UTC; see the text).}\label{tab:stages}
\begin{tabular}{@{}l l p{3.6cm}@{}}
\toprule
Stage & Jobs & What it measures\\
\midrule
Stage 0 & 17060165 & smoke test\\
Stage C & 17062103 & closed $\mu(n)$, 300k pool, FCFS, $n\in\{16,20,24,32\}$ (node gpu-1)\\
 & 17062104 & FCFS, $n\in\{48,96,300\}$ (gpu-3)\\
 & 17062105 & LPM, $n\in\{48,300\}$ (gpu-2)\\
E0a & 17072634--37 & $\lmiss$, cache off\\
E0b & 17084750--55 & from-empty stability\\
E2 & 17093387--404 & registered paired test, 2 rates $\times$ 3 seeds $\times$ 3 arms\\
native C & 17096856 & $\mu(n)$, native pool, $n\in\{24,48,64,80\}$ (gpu-2)\\
 & 17096857 & $n\in\{96,128,300\}$ (gpu-1)\\
native E0a & 17096858 & $\lmiss$, native pool\\
grid & 17097086--103 & BURST arm, 3 policy cells\\
vLLM & 17097972--74 & install, GPU check, smoke\\
native E0b & 17100938--39 & from-empty stability, native pool\\
vLLM E0a & 17100940 & $\lmiss$, vLLM, cache off\\
native E2 & 17102173--90 & registered paired test, native pool, 2 rates $\times$ 3 seeds $\times$ 3 arms\\
vLLM BURST & 17102318--23 & BURST arm, vLLM, exploratory\\
vLLM C & 17162959 & $\mu(n)$, vLLM, 300k pool, $n\in\{16,20,24\}$ (gpu-3)\\
 & 17162960 & $n\in\{28,32,40\}$ (gpu-3)\\
 & 17162961 & $n\in\{48,96,300\}$ (gpu-2)\\
vLLM E0b & 17162962--63 & from-empty stability, vLLM\\
vLLM E2 & 17166121--38 & registered paired test, vLLM, 2 rates $\times$ 3 seeds $\times$ 3 arms\\
\midrule
\multicolumn{3}{@{}l}{\emph{Rounds 11 to 14 (Tables~\ref{tab:frozen}--\ref{tab:frozen3})}}\\
H200 & 17166807 & smoke, H200 NVL, 300k pool\\
 & 17166810--13 & E0a and its 64-session check\\
 & 17166818--22 & Stage C, $n=16$--300 (nine runs, with the next row)\\
 & 17166828, 30, 32, 34 & Stage C, the rest of the nine runs\\
 & 17168377--82 & E0b\\
 & 17169668--71 & E0b reseeds, added after 17168377--82 were read (post hoc)\\
 & 17169612--14 & Stage C, $n=36$, 44, 52\\
 & 17169672--80, 82--87 & Stage C, E0b, 600k pool\\
 & 17171020--37 & E2, 600k pool\\
Llama & 17166808 & smoke, Llama-3.1-8B\\
 & 17166814--16 & Stage C\\
 & 17166823, 29, 31, 33 & E0a and its 64-session check\\
 & 17169565--70 & E0b\\
 & 17170409--26 & E2\\
pools & 17168864--81 & Stage C, E0b; 150k, 600k\\
BASIN-1 & 17168500--05 & bursts of 100, then 25\\
E0F & 17168506--07 & from-empty, seeds 3--5\\
LPM-INF & 17168743--45 & BURST, no LPM fallback\\
BASIN-2 & 17169893--901 & bursts of 40, 60, then 50\\
TOOLSLOW & 17169902--07 & $\delta=10$, then 30\,s\\
SHED & 17169908--11 & $\tau=600$, then 3600\,s\\
rescue & 17168965 & smoke test of the gate (never decisive)\\
 & 17168966--71 & R1, 300k pool, 2 rates $\times$ 3 seeds\\
 & 17168972--74 & R4, native pool, 3 seeds\\
\bottomrule
\end{tabular}
\end{table}

\noindent The accounting export \fpath{provenance/sacct_cited_jobs_2026-11-01.txt} was made after the last cited job had ended; its times are Slurm's local time, CDT until 02:00 on 2026-11-01 and CST after it, and \fpath{provenance/sacct_cited_jobs_2026-11-01_utc.txt} gives the same records in UTC. The earlier export \fpath{provenance/sacct_cited_jobs_2026-10-30.txt} covers every job through the vLLM E0b, and the data directories of the vLLM paired test and of the BASIN-1, fresh-seed, rescue, pool-sweep, BASIN-2, TOOLSLOW and SHED jobs hold their own copies (\fpath{sacct_17166121-17166138.txt}, \fpath{sacct_17168500-17168507.txt}, \fpath{sacct_1716896x.txt}, \fpath{sacct_17168864-17168881.txt}, \fpath{sacct_17169893-17169901.txt}, \fpath{sacct_17169902-17169911.txt}).

\begin{table}[htbp]
\centering\footnotesize
\caption{Workload and engines of the real-engine runs.}\label{tab:workload}
\begin{tabular}{@{}l l@{}}
\toprule
First prompt $L_0$ & 2000 tokens\\
Tool output per turn $\Delta$ & 550 tokens\\
Output per turn $d$ & 100 tokens, greedy, \texttt{ignore\_eos}\\
Turns $T$ & $\min(\mathrm{Geometric}(1/15),60)$, $E[T]=14.761$\\
Tool time $Z$ & exponential, mean 5\,s\\
Prompt at turn $k$ & $2000+650(k-1)$ tokens\\
Output feedback & server-returned token ids\\
\midrule
Engine & SGLang 0.5.22, FlashInfer; FA3 on the H200\\
Scheduling & FCFS, waiting-prefix refresh on\\
Batching & chunked prefill 8192, $\le256$ running\\
KV pool & 300,000 tokens; native 955,757\\
 & pool sweep 150,000, 600,000; H200 also 600,000\\
Context limit & 49,152 tokens\\
Model & Qwen2.5-7B-Instruct, rev.\ a09a3545\\
 & one cell: Llama-3.1-8B-Instruct, rev.\ d10aef7\\
 & (NousResearch mirror, shards identical to Meta's)\\
GPU & A100 80GB PCIe, 81,920\,MiB, driver 580.178.04\\
 & two cells: H200 NVL, 143,771\,MiB, same driver\\
\midrule
vLLM & 0.32.0, FCFS (default), prefix caching on\\
 & chunked prefill 8192, $\le256$ running\\
 & 300,000 tokens (18,750 blocks of 16)\\
\bottomrule
\end{tabular}
\end{table}

Table~\ref{tab:workload} lists the workload and engine settings of the real-engine runs.

\paragraph{The sequence of stages.} \emph{Stage~0}, a smoke test of server and client on one GPU, came first; its server logs and those of Stage~C calibrated the simulator (\S\ref{sec:sim}). After E2 we repeated the sequence with the native pool (Stage~C, E0a, E0b, a second registration and \emph{native E2}) and ran an exploratory policy grid. On vLLM we first measured $\lmiss$ and ran the BURST arm alone under an exploratory registration, and later ran its own Stage~C and E0b, a third registration and its paired test (\S\ref{sec:vllm}). An exploratory replication of E0b later used seeds 3--5 at two rates (\S\ref{sec:edge}). Later cells repeat the first stages on an H200 at the same pool and at 600k tokens (\S\ref{sec:h200}), with Llama-3.1-8B, and on the A100 with pools of 150k and 600k tokens; the Llama-3.1-8B and 600k H200 cells have registered paired tests, and registered tests of a rescue (\S\ref{sec:rescue}), of other burst sizes (\S\ref{sec:basin}), of a pause in arrivals and of a tool slowdown follow the paired-test protocol; \S\ref{sec:frozen} lists their frozen predictions and the outcomes in hand.

\paragraph{Stage 0.} The smoke test (job 17060165) verified that every engine flag exists; that cached tokens are reported (a repeated 64-token prompt had 63 tokens cached); that the server returns output token ids; that contexts up to 40,350 tokens are accepted; that the KV pool is exactly 300,000 tokens; that client and server agree on prompt length (0 mismatches over 3191 requests); and that the flush empties the cache (canary 1023 cached tokens before, 0 after).

\paragraph{Stage C and E0a.} Closed population; all sessions start at $t=0$ and a finished session is replaced at once; 900\,s warmup and an 1800\,s measured window; one run per $n$ (seed 0); flush and canary check between runs. The native-pool and vLLM closed curves used the same protocol, with the populations of Table~\ref{tab:stages}. E0a disables the radix cache and runs 32 sessions with seeds 0, 1 and 2, plus 64 sessions with seed 0 as a saturation check. Its per-run values over the window were (sessions/s, turns-based sessions/s, uncached tokens/s): seed 0, 0.0567, 0.0643, 11,076; seed 1, 0.0817, 0.0773, 11,342; seed 2, 0.0767, 0.0730, 11,312; 64 sessions, 0.0544, 0.0676, 11,431. The spread comes from in-flight composition: the mean turn index of the requests in the window was 15.9, 13.2 and 14.1 for seeds 0, 1 and 2. vLLM E0a ran the same protocol with vLLM's prefix caching off: seed 0, 0.0539, 0.0603, 10,208; seed 1, 0.0728, 0.0705, 10,487; seed 2, 0.0700, 0.0662, 10,352; 64 sessions, 0.0467, 0.0631, 10,531.

\paragraph{E0b.} Poisson session arrivals from an empty cache and queue at 1.25, 1.5, 1.75, 2.0, 2.5 and 3.0 $\lref$; one 3000\,s run per rate; stable if the final queue is below 50 and the fitted rise of the queue over the last 1500\,s is at most 10. All runs used arrival seed 0, so they share one random-number stream and the arrivals at each rate are a subset of those at the next higher rate. The six runs received 302, 359, 416, 472, 585 and 697 sessions, 8.3--12.6\% above $\lambda T$ and about 2.1 Poisson standard deviations above it in every run. Native E0b used the same rule at 1.5, 2.5, 3.5, 4.5, 5.5 and 6.5 times the native $\lref$ (361, 587, 811, 1034, 1257 and 1478 sessions, 5.6--11.7\% above $\lambda T$), and vLLM E0b the multiples of the 300k runs of vLLM's $\lref$ (279, 331, 383, 435, 539 and 642 sessions, 8.6--13.3\% above $\lambda T$; stable through 2.0, final queue 81.6 and 247.1 at 2.5 and 3.0).

The fresh-seed E0b (E0F; jobs 17168506--07) repeated 2.0 and 2.5 $\lref$ on the 300k pool with arrival seeds 3, 4 and 5 under the same rule. All six runs were stable, with final queue 0, so the frozen replication rule, which needed 2.0 $\lref$ stable in all three seeds and 2.5 $\lref$ unstable in at least two, is not met. The nominal bracket of 2.0 to 2.5 $\lref$ therefore comes from the single seed-0 grid: over seeds 0 and 3--5, 2.0 $\lref$ is stable in 4 of 4 runs and 2.5 $\lref$ in 3 of 4. Measured instead in realized arrival rates (base arrivals in the 3000\,s run over $3000\,\mathrm{s}\times\lref$), the twelve from-empty runs of seeds 0 and 3--5 are stable up to 2.665 $\lref$ and unstable from 2.726 $\lref$, monotone in the realized rate, so the edge lies between 2.665 and 2.726 $\lref$. The seed-0 run at 2.5 $\lref$ is the one at 2.726: it received 9\% more arrivals than its nominal rate. Either way, the registered rates of 1.35 and 1.65 $\lref$ lie inside the window. The simulator's blind prediction was stable in all six runs; the heuristic H2 of \S\ref{sec:fp-heur} gave each 2.5 $\lref$ run a probability of 0.69 of being unstable and was too pessimistic. This test is exploratory: it has no registration of its own and does not change the E2 registration, whose window was fixed by the seed-0 bracket.

\paragraph{E2.} One arm per 8700\,s job, each on a fresh server; burst of 200 fresh sessions within 60\,s at $t_b=1200$\,s for BURST and CAPPED; CAPPED gate $\ncap=20$; evaluation window $[t_b+4500,\,8700]$\,s. The registered rule (Table~\ref{tab:rule}) was applied by the analysis tool frozen at the registration commit, with the same parameter overrides that the jobs themselves used. Native E2 used the same protocol with the native registration's rates, reference rates and gate $\ncap=64$. Twelve of its 18 runs started after the cluster checkout moved from the registration commit to \texttt{a78ce07}, which adds only the vLLM registration, its input file and two notes files (the research log and the status notes), with no code change, and its verdict ran on that commit; two of its 18 servers sized the pool one token smaller, at 955,756. The vLLM paired test uses the same protocol with the vLLM registration and $\ncap=20$. It uses the same seeds, rates and reference rate as the exploratory vLLM BURST runs, so its BURST arm repeats their workloads rather than replicating them independently.

\paragraph{Why the rule uses tokens.} Criterion (a) uses tokens because a trapped engine with a growing backlog serves mostly young sessions, whose turns are short, so at the all-miss token rate it completes more turns per second than the stationary all-miss mix, and sessions faster than $\lmiss$ (\S\ref{sec:fluid}): on traps simulated in our earlier simulator (\S\ref{sec:sim}), a session-rate test fired on at most 1 of 29 at 2 $\lmiss$ and none of 39 at 2.5, while the token test fired on all 84 traps at 1.5--2.5 $\lmiss$.

\paragraph{Drain-time estimates.} The registrations' feasibility check used $b/(\mu(c)-\lambda)$, which ignores the $c-n_1$ free slots that a burst fills at once: 1571\,s at 1.35 and 1889\,s at 1.65 $\lref$ for E2, 668\,s and 1043\,s at 2.5 and 4.0 $\lref$ for the native test, and 1529\,s and 1800\,s for vLLM. With the populations present at the burst (5--12 sessions at 300k, 11--32 with the native pool, 5--11 on vLLM), the fluid estimate $(n_1+b-c)/(\mu(c)-\lambda)$ of \S\ref{sec:props} is 1453--1508\,s and 1757--1804\,s, 491--518\,s and 840--876\,s, and 1414--1460\,s and 1674--1692\,s. The CAPPED gate's queue emptied 1547--1904\,s and 2021--2693\,s after the burst on E2, 3--52\% longer than this estimate; 512--583\,s and 837--937\,s with the native pool, between 4\% shorter and 13\% longer; and 1521--2735\,s on vLLM, 1.05--1.63 times it.

\paragraph{Rescue (R1, R4).} One rescue run per 8700\,s job, each on a fresh server, from a frozen clone at the registration tag; the client's arrival plan is the BURST plan of the same seed and rate, byte for byte, so until $t_{\mathrm{on}}$ the rescue run replays the BURST run's arrival plan (a fresh run on a fresh server, not the same run). At $t_{\mathrm{on}}$ the gate starts to bind with cap $c$ by attrition and pause: a session whose tool call ends while more than $c$ sessions are admitted gives up its slot and joins a resume queue, which is served before the queue of new arrivals, each in FIFO order; a paused session's KV is not protected while it waits. The cap is removed at the first 10\,s tick at or after $t_{\mathrm{on}}+600$\,s at which both queues have been empty for 300\,s; there is no forced release. The registered per-run rule uses the follow-up hit fraction, the share of a session's turns after its first, attributed by completion time, with more than half of the prompt cached: (a) VALID if the follow-up hit fraction over $[t_{\mathrm{on}}-300,\,t_{\mathrm{on}})$ is below 0.2 and more than 50 requests or sessions are outstanding at $t_{\mathrm{on}}$ (requests in flight plus sessions waiting at the gate or paused), on the registered hardware; (b) RELEASED by 7200\,s; (c) POST-GOOD over $[t_{\mathrm{rel}}+300,\,8700)$: every 300\,s window, tiled from the start with a short remainder covered by the window ending at 8700\,s, has a follow-up hit fraction of at least 0.5, the whole interval above 0.8, an uncached token rate below 0.5 times the all-miss reference, and the registered RECOVERED test holds over the last 300\,s; a window with no completed follow-up fails (c) without being a relapse. RESCUED is (a), (b) and (c); RELAPSE is any such window below 0.2. A same-seed BURST control, from the paired test, is STUCK if its follow-up hit fraction over $[t_{\mathrm{on}},\,8700)$ is below 0.2 and it is not RECOVERED over the last 300\,s. A cell passes if at least 2 of its 3 seeds are RESCUED and every same-seed control is STUCK; R1 is refuted if 2 of its 6 runs relapse or are not released by 7200\,s, and R4 if 2 of its 3 relapse; a refutation takes precedence over a pass, and a failed session in any counted run withdraws the claim that nothing was shed. The pairing requires common random numbers with the committed BURST and NOBURST runs (the BURST plan byte-identical, the NOBURST plan the rescue plan without the burst sessions), equal resolved rates, references, hardware and burst time. Each job ran only a mechanical check of the gate (\texttt{analysis.rescue check-run}: engage within 10\,s of $t_{\mathrm{on}}$, release on the 10\,s grid after the 300\,s hold, no request errors), which passed in all ten jobs; no hit fraction of a rescue run was read before all runs of its cell had ended. The controls ran on 2026-10-24 to 26 as part of the paired tests, not alongside the rescue runs. A smoke job under its own registration checked the gate's mechanics before the registered runs and is never decisive.

\begin{table}[htbp]
\centering\footnotesize
\caption{The registered decision rule.}\label{tab:rule}
\begin{tabular}{@{}l l@{}}
\toprule
Rate estimator & uncached prefill tokens/s\\
Reference rate & E0a all-miss, 11,243 tokens/s\\
 & \quad (native: 11,282; vLLM: 10,349)\\
$\lmiss$ in the rule & the cell's $\lref$\\
Rate tolerance & $\pm0.15$\\
Backlog slope & least squares over 10\,s bin means of\\
 & \quad sessions in the system, gate included\\
Minimum window & 600\,s\\
Trapped & $\lambda>\lmiss$, rate within tolerance,\\
 & \quad and slope $>0.5\,(\lambda-\lmiss)$\\
Recovered & hit $>0.8$ and outstanding $<20$, last 300\,s\\
Draining & slope $<0$\\
Indeterminate & none of the above\\
Precedence & trapped, recovered, draining, indeterminate\\
Hit metric & turns with cached/prompt $>0.5$\\
Pass & $\ge2/3$ of $\ge3$ seeds with all three\\
 & \quad arm conditions met in the seed\\
\bottomrule
\end{tabular}
\end{table}

\paragraph{Backlog definitions.} The figures and Table~\ref{tab:e2} report the backlog as sessions in the system (arrived minus ended), which includes sessions asleep in tool calls (8--12 in the good state at 300k, 16--34 with the native pool). The registered RECOVERED test uses the outstanding backlog (requests in flight plus sessions held at the gate). The simulator's predictions count waiting requests plus the gate queue. The rescue rule's outstanding work also counts sessions paused by the gate (Fig.~\ref{fig:rescue}).

\section{The registrations}\label{app:prereg}
\begingroup\raggedright
The registration was filled by formula from committed copies of the Stage~C curve, the E0a calibration and the E0b bracket. The window was $(1,\,2.0)\,\lref$, its upper edge E0b's largest all-stable multiple (below $\mu^*/\lref=3.26$), a nominal value from E0b's single seed-0 grid; fresh seeds later put the from-empty edge, in realized arrival rates, between 2.665 and 2.726 $\lref$, and both registered points lie inside the stable range on either reading (\S\ref{sec:edge}). The rates followed a rule proposed in the runbook and adopted in the project status notes on 2026-10-22, the day before E0b was submitted (Appendix~\ref{app:prereg}): one third and two thirds of the way across the window, rounded to 0.05, which gives 1.35 and 1.65 $\lref$ (0.0966 and 0.1180\sps). The cap was the smallest $n$ with $\mu(n)\ge0.95\mu^*$: $\ncap=20$, since $\mu(20)=0.224$ is about 4\% below $\mu(24)=0.233$, inside the rule's 5\%. The file was committed and reached the project's GitHub repository at 15:09:54 CDT on 2026-10-24, before the E2 jobs were submitted at 15:13 (Table~\ref{tab:timeline}). We call a test pre-registered when its registration file was committed and pushed to GitHub before its runs; this is not an entry in a third-party registry. Registration and the frozen predictions of \S\ref{sec:sim} are rigor, not a contribution: they let a reader check that each rule and each prediction existed before its outcome. The design follows the model-to-system pipeline of Isaacs et al., in which an abstract model chooses the experiments that then test it~\citep{isaacs2025analyzing}, and pre-registered gates have been used in ML-systems measurement before~\citep{suram2026cacheable}. Appendix~\ref{app:prereg} lists every registration we filled, with its status, and Appendix~\ref{app:timeline} every prediction round with its time record; the files carry no timestamp other than GitHub's push log. The repository we publish is a redacted copy of the development repository whose \fpath{provenance/} directory ties every registration file, by hash chains, to the development repository's push log (availability statement). Commit times are set by the committer, so the external record is GitHub's push log, archived as \fpath{provenance/github_activity_2026-10.json} (Appendix~\ref{app:timeline}). The job script refuses to start unless the file is committed and unmodified and the job's cell, rates, cap, reference rates, hardware and protocol match it, and writes its SHA-256 into every run; the analysis calls a verdict decisive only for the registered cell and rates with every run bound to the file and started after its commit (Appendix~\ref{app:prereg}); the label records this binding, not statistical strength, and the analysis checks commit order, not push time, so its own label ``pre-registered'' does not by itself meet our definition. The native registration was filled and enforced the same way and reached GitHub 10 minutes before its jobs were submitted (\S\ref{sec:e2native}). The vLLM registration was filled from vLLM's own Stage~C curve, E0a calibration and E0b bracket, and does not meet the definition: it was committed 37\,s before its 18 jobs were submitted, and the cluster's record of the submission names the file but not its hash, but it reached GitHub 2\,min after the submission, when two of the jobs had started and none could have ended, so we call this test registered, not pre-registered (\S\ref{sec:vllm}). On the H200 with the 300k pool the rule could not be applied: its window was $(1,\,1.3)\,\lref$ (the largest all-stable E0b multiple, below $\mu^*/\lref=1.99$), and the points the rule gives, 1.10 and 1.20 $\lref$, lie outside the band $(1.1,\,1.17)\,\lref$ that its 10\% margin at each edge of the window leaves, so the fill returned INCOMPLETE and that cell has no registered test (\S\ref{sec:frozen}). With the 600k pool the window was $(1,\,2.5)\,\lref$ and the rule filled a registration (Table~\ref{tab:ledger}).

The registration is \fpath{configs/agentload_prereg_2026-10-24.yaml}, status FILLED, committed in \texttt{d279b7b} at 15:09:12 CDT on 2026-10-24 and pushed to GitHub at 15:09:54, before the E2 jobs were submitted at 15:13:39, SHA-256 \texttt{2a6d680a9615bba962a5673d6db6b3e4\allowbreak ca5102bf778ca974cea46147f5cc0ac6}. Its fields:

\begin{itemize}
\item \textbf{Cell and inputs.} Engine cell \texttt{sglang-fcfs-r1} (SGLang FCFS, waiting-prefix refresh on); committed copies of the Stage~C curve, the E0a calibration and the E0b bracket under \fpath{prereg/inputs/2026-10-23/}; protocol 8700\,s, burst of 200 within 60\,s at 1200\,s.
\item \textbf{Measured.} $\mu(n)$ in sessions/s: 0.1756 ($n=16$), 0.2239 (20), 0.2333 (24), 0.1150 (32), 0.0489 (48), 0.0611 (96), 0.1078 (300). $\mu^*=0.2333$; $\ncap=20$ by the smallest-within-5\% rule (argmax 24; flagged ambiguous). $\lmiss=0.07167$ (sessions) and 0.07153 (turns-based, $=\lref$); all-miss rate 11,243 tokens/s; E0b largest all-stable multiple 2.0, first unstable 2.5, monotone (arrival seed 0 only; the fresh seeds of Appendix~\ref{app:protocol} came later); hardware A100 80GB PCIe.
\item \textbf{Predicted boundaries.} Recovery threshold $\lref$ (not $\mu(300)$, which a synchronized start biases upward). Window $(0.0715,\,0.1431)$\sps{} $=(1.0,\,2.0)\,\lref$, upper edge from E0b. CAPPED drain must fit within 80\% of the 7200\,s available after the burst (5760\,s), which holds for $\lambda$ up to 2.64 $\lref$; predicted drain 1571\,s at 1.35 and 1889\,s at 1.65.
\item \textbf{Rates.} 1.35 and 1.65 $\lref$ by the one-third/two-thirds rule, rounded to 0.05, inside the margin band $(1.1,\,1.8)$; seeds 0, 1, 2. On 2026-10-22 the runbook proposed the rate rule and a cap of 20, with keeping 24 as the alternative (\texttt{288b801}, 13:13 CDT); both rules were adopted in the project status notes at 13:50 (\texttt{4435c51}) and reached GitHub at 17:45, before E0b was submitted on 2026-10-23; the registration's own text dates them 2026-10-23.
\item \textbf{Predictions.} Inside the window: BURST trapped, NOBURST good through the run horizon, CAPPED recovered if the drain fits. Below $\lref$: BURST not trapped. Above the window: NOBURST may escape, so no such rate was registered.
\item \textbf{Kill conditions.} NOBURST not good in at least one third of the seeds at a registered rate; BURST not trapped in two thirds of the seeds at every registered rate; CAPPED not recovering while BURST traps and the drain fits; a window that is empty or narrower than the noise.
\item \textbf{Warnings recorded at filling.} $|\mu(300)/\lmiss-1|=0.507$ (a check only), and the ambiguous cap.
\end{itemize}

\noindent The verdict is in \fpath{results/agentload_analysis_burst2/20261025-144012/} inside the E2 data directory: the registered analysis passes at 1.35 and 1.65 $\lref$ (the tool's labels, assigned by commit order, are pre-registered and decisive; by the definition of \S\ref{sec:prereg} the test is pre-registered). The three grid cells have their own derived registrations, \fpath{configs/agentload_prereg_2026-10-24_grid_<cell>.yaml} for the cells \texttt{sglang-fcfs-r0}, \texttt{sglang-lpm-r1} and \texttt{sglang-lpm-r0}, with status EXPLORATORY; the analysis never calls them decisive.

\paragraph{The native registration.} \fpath{configs/agentload_prereg_2026-10-25_native.yaml}, status FILLED, committed in \texttt{b8c19ec} at 17:31:16 CDT on 2026-10-25 and pushed to GitHub at 17:31:56, 10 minutes before the native E2 jobs were submitted (17:42:29), SHA-256 \texttt{21bf7c03e0630b65b80d086058815e12\allowbreak c5ba9fbe41e9864d3dd0183b9bf776c1}; the same template, formulas, trap rule, kill conditions and protocol as above. Cell \texttt{sglang-fcfs-r1-native}; inputs under \fpath{prereg/inputs/2026-10-25-native/}. Measured: $\mu(n)$ 0.2650 ($n=24$), 0.4311 (48), 0.4789 (64), 0.4772 (80), 0.1228 (96), 0.0667 (128), 0.1006 (300); $\mu^*=0.4789$; $\ncap=64$ by the smallest-within-5\% rule (also the argmax; flagged ambiguous, $n=80$ within 0.35\%); $\lref=0.07180$ (turns-based); all-miss rate 11,282 tokens/s; native E0b largest all-stable multiple 5.5, first unstable 6.5, monotone. Window $(0.0718,\,0.3949)$\sps{} $=(1,\,5.5)\,\lref$, upper edge from native E0b, below $\mu^*/\lref=6.67$. Rates 2.5 and 4.0 $\lref$ (0.1795 and 0.2872\sps) by the same rule, inside the margin band $(1.1,\,4.95)$; the rule was recorded in the project status notes (\texttt{6cb43e9}, pushed 15:08 CDT) before native E0b was read. Predicted drains 668\,s and 1043\,s; the drain fits up to 6.19 $\lref$. Warnings recorded at filling: $|\mu(300)/\lmiss-1|=0.401$, the ambiguous cap, and an incomplete grid (no $n=16$, 20 or 32). The verdict is \fpath{results/agentload_analysis_burst2/20261026-100653/} inside the native E2 data directory: the registered analysis passes at 2.5 and 4.0 $\lref$ (with the same tool labels; pre-registered by the definition of \S\ref{sec:prereg}).

\paragraph{The exploratory vLLM registration.} \fpath{configs/agentload_prereg_2026-10-24_grid_vllm-fcfs.yaml}, status EXPLORATORY, derived from the E2 registration, committed in \texttt{a78ce07} at 18:06:18 CDT on 2026-10-25 and pushed to GitHub at 18:06:33, 46\,s before the vLLM BURST jobs were submitted, SHA-256 \texttt{c5037cac6003c4e98da89f88dae953bf\allowbreak e68bb21a27a4e3e5e3dab3c8891aeab5}. Cell \texttt{vllm-fcfs}, BURST arm only; the window, multiples, protocol and trap rule are inherited from \texttt{sglang-fcfs-r1}, and the reference rates come from \fpath{prereg/inputs/2026-10-25-vllm/calibration.json}, a byte-identical copy of the vLLM E0a calibration. Its predicted-boundary block is copied unchanged from the E2 registration (window 0.0715--0.1431\sps, all-miss rate 11,243 tokens/s), but the runs used vLLM's references, 0.06566\sps{} and 10,349 tokens/s (\texttt{lambda\_ref\_resolved} and \texttt{uncached\_tokens\_ref\_resolved} in each run's \fpath{meta.json}), so their rates were 1.35 and 1.65 times vLLM's $\lref$. The verdict, \fpath{results/agentload_analysis_burst2/20261026-100701/} inside the native E2 and vLLM BURST data directory, reports per-run states and no registered verdict.

\paragraph{The vLLM registration.} \fpath{configs/agentload_prereg_2026-10-30_vllm.yaml}, status FILLED, SHA-256 \texttt{8f2d5b44d6365795fedde8bf673d51e5\allowbreak 4bf8a685a14d563c86de44492bb95f84}; the same template, formulas, trap rule, kill conditions and protocol as the E2 registration. Cell \texttt{vllm-fcfs}; inputs: the vLLM Stage~C curve and E0b bracket under \fpath{prereg/inputs/2026-10-30-vllm/}, byte-identical to the analyses in the vLLM Stage~C and E0b data directories, and the vLLM E0a calibration above. Measured: $\mu(n)$ 0.1750 ($n=16$), 0.2194 (20), 0.2183 (24), 0.1189 (28), 0.0433 (32), 0.0411 (40), 0.0439 (48), 0.0572 (96), 0.1050 (300); $\mu^*=0.2194$; $\ncap=20$ by the smallest-within-5\% rule (also the argmax; flagged ambiguous, $n=24$ within 0.5\%); $\lref=0.06566$ (turns-based); all-miss rate 10,349 tokens/s; vLLM E0b largest all-stable multiple 2.0, first unstable 2.5, monotone. Window $(0.0657,\,0.1313)$\sps{} $=(1,\,2.0)\,\lref$, upper edge from vLLM E0b, below $\mu^*/\lref=3.34$. Rates 1.35 and 1.65 $\lref$ (0.0886 and 0.1083\sps) by the same rule, inside the margin band $(1.1,\,1.8)$; this rule and the cap rule were recorded in the project status notes (\texttt{836063a}, pushed 12:16 CDT on 2026-10-30) after the vLLM Stage~C and E0b jobs were submitted and before any of them had ended. Predicted drains 1529\,s and 1800\,s; the drain fits up to 2.81 $\lref$. Warnings recorded at filling: $|\mu(300)/\lmiss-1|=0.599$ and the ambiguous cap. The file was filled from five of the six vLLM E0b runs while the 3.0 $\lref$ run was still running; with 2.5 $\lref$ unstable that run could not change the edge, and it was unstable too, so the six-run bracket equals the registered one. The file was committed in \texttt{5a6fbcd} at 16:23:59 CDT, 37\,s before the 18 jobs of the paired test (17166121--38) were submitted at 16:24:36 (the research log gives 16:30), and the jobs run from the cluster's copy of the repository at that commit; Slurm's record of the submission names the file but not its hash. It reached GitHub at 16:26:35, 1\,min 59\,s after the submission, when two of the jobs had started (at 16:25:42 and 16:26:17) and none could have ended. This registration therefore does not meet the definition of \S\ref{sec:prereg}: it was committed before the jobs were submitted, but pushed after the first two had started, so we call this test registered, not pre-registered. As for E2, the job script writes the file's SHA-256 into every run's \fpath{meta.json}.

The verdict, \fpath{results/agentload_analysis_burst2/20261031-020837/} in the paired test's data directory (\fpath{results/nova/e2_vllm_1716612x_pulled_20261031-020956/}), reports a pass at both rates; all 18 runs' \fpath{meta.json} files carry the registration's SHA-256. The tool also prints the labels pre-registered and decisive, which it assigns by commit order; by the definition of \S\ref{sec:prereg} the test is registered, not pre-registered. The frozen mechanism predictions for this test (\fpath{configs/r11_mechanism_predict.yaml}, \fpath{notes/PREDICTIONS_r11_mechanism.md}; commit \texttt{f1338e3}, pushed at 18:20:13 CDT on 2026-10-30, when three of the 18 jobs were running and none had ended; the first ended at 18:55:07) were evaluated with the frozen scripts, run unchanged. Byte copies of the frozen files and their outputs are under \fpath{r11_mechanism/} in the same directory, and the scored decision is \fpath{derived/r11_mechanism_decision.json}, which records two detector false alarms in CAPPED runs, one more than the frozen rule allows (\S\ref{sec:vllm}).

\paragraph{The rescue registrations.} \fpath{configs/agentload_prereg_2026-10-31_rescue.yaml} (R1, cell \texttt{sglang-fcfs-r1-rescue}) and \fpath{configs/agentload_prereg_2026-10-31_rescue_native.yaml} (R4, cell \texttt{sglang-fcfs-r1-rescue-native}), status FILLED, SHA-256 \texttt{502dee9e3107952a2f968c1b59b8ee3d\allowbreak 4063870d72f0dfac1261e2b2e2316691} and \texttt{ddd776f822755ba50472668b4af7dba3\allowbreak 538354d98feb2579b8cca422b5215bb1}, were committed in \texttt{f0bc3ad} at 22:10:47 CDT on 2026-10-30 (tag \texttt{r11-rescue-reg}) and reached GitHub at 22:39:02, 22 minutes before the ten jobs (a smoke job and the nine registered runs, 17168965--74) were submitted at 23:01:12--13; by the definition of \S\ref{sec:prereg} both tests are pre-registered. They reuse the inputs, rates and caps of the E2 and native E2 registrations, under their own cell labels, and add the rescue protocol, the per-run rule, the decision rule and the refutation conditions of Appendix~\ref{app:protocol}, with the rule's constants frozen in \fpath{analysis/rescue.py} at the tag. Three changes were accepted after a review of the simulator's predictions and before any rescue run: a refutation takes precedence over a pass, the no-shed part of a pass is withheld if any counted run has failed sessions, and R4, like R1, requires stuck same-seed BURST controls. R4's engage time of 2100\,s replaced 2400\,s after the simulator had predicted 13 of 20 rescued at 2400\,s (\fpath{notes/PREDICTIONS_r11_rescue.md}). The verdicts, \fpath{results/agentload_analysis_rescue_r1/20261031-071720/} and \fpath{results/agentload_analysis_rescue_r4/20261031-071748/} in the data directory \fpath{results/nova/rescue_1716896x_pulled_20261031-070502/}, were computed on a clean tree with the registrations and the simulator's predictions as inputs: each reports a pass, with no rejected run and no pairing, protocol or registration problem, and every run's \fpath{meta.json} carries its registration's SHA-256. The tool labels each verdict decisive, a label it assigns by commit order.

\paragraph{Ledger of registrations.} Table~\ref{tab:ledger} lists every registration file we filled, with its status. Commit times are the committer's; push times come from GitHub's push log (\fpath{provenance/push_times_2026-11-01.tsv}, built from the export of 2026-11-01, or the archived log itself for the grid cells). Each is the first branch update that contains the commit; the BASIN-1 push also carried \texttt{9b67c26}, the commit of tag \texttt{r11-basin-reg}.
\par\endgroup

\begin{table*}[tbp]
\centering\footnotesize
\caption{Every registration, in order. FILLED registrations bind a test's cell, rates, cap, references, hardware and protocol, and the job script refuses to run a test without one; EXPLORATORY ones are derived from a FILLED parent and are never called decisive. Times CDT.}\label{tab:ledger}
\begin{tabular}{@{}>{\raggedright\arraybackslash}p{5.0cm} l l >{\raggedright\arraybackslash}p{2.2cm} >{\raggedright\arraybackslash}p{4.0cm}@{}}
\toprule
File (\fpath{configs/agentload_prereg_*}) & status & committed & on GitHub & runs and outcome\\
\midrule
\fpath{2026-10-24.yaml} (E2) & FILLED & \texttt{d279b7b} 10-24 15:09:12 & 15:09:54 & 18 runs; pass (\S\ref{sec:e2})\\
\fpath{2026-10-24_grid_sglang-*.yaml} (3 grid cells) & EXPLORATORY & \texttt{499c0df} 10-24 23:30:55 & 23:31:32 & 18 BURST runs (\S\ref{sec:grid})\\
\fpath{2026-10-25_native.yaml} (native E2) & FILLED & \texttt{b8c19ec} 10-25 17:31:16 & 17:31:56 & 18 runs; pass (\S\ref{sec:e2native})\\
\fpath{2026-10-24_grid_vllm-fcfs.yaml} & EXPLORATORY & \texttt{a78ce07} 10-25 18:06:18 & 18:06:33 & 6 BURST runs (\S\ref{sec:vllm})\\
\fpath{2026-10-30_vllm.yaml} (vLLM E2) & FILLED & \texttt{5a6fbcd} 10-30 16:23:59 & 16:26:35, after two jobs started & 18 runs; pass; registered, not pre-registered\\
\fpath{2026-10-30_b100.yaml}, \fpath{_b25.yaml} (BASIN-1) & FILLED & \texttt{9128995} 10-30 19:51:20 & 20:04:30 & 6 BURST runs; pass (Table~\ref{tab:frozen2})\\
\fpath{2026-10-30_grid_sglang-lpminf.yaml} & EXPLORATORY & \texttt{5064580} 10-30 21:02:19 & 21:02:46 & 3 BURST runs; no collapse on the two informative seeds, no recovery (Table~\ref{tab:frozen3})\\
\fpath{2026-10-31_rescue.yaml} (R1), \fpath{_rescue_native.yaml} (R4), \fpath{_rescue_smoke.yaml} & FILLED & \texttt{f0bc3ad} 10-30 22:10:47 & 22:39:02 & 9 runs, submitted 23:01; R1 and R4 pass (\S\ref{sec:rescue}); the smoke job is never decisive\\
\fpath{2026-10-31_b40.yaml}, \fpath{_b50.yaml}, \fpath{_b60.yaml} (BASIN-2) & FILLED & \texttt{621d88f} 10-31 01:58:38 & 02:10:48 & 9 runs, submitted 04:21; bracketed: 40 good, 60 stuck, 3 of 3 each (\S\ref{sec:basin})\\
\fpath{2026-10-31_toolslow_d10.yaml}, \fpath{_d30.yaml} & FILLED & \texttt{b006291} 10-31 02:27:40 & 02:28:09 & 6 runs, submitted 04:21; rule not met: 10\,s good in 3 of 3, 30\,s stuck in 1 of 3 (\S\ref{sec:frozen})\\
\fpath{2026-10-31_shed_t600.yaml}, \fpath{_t3600.yaml} & FILLED & \texttt{b006291} 10-31 02:27:40 & 02:28:09 & 4 runs, submitted 04:21; pass: 600\,s stuck, 3600\,s good, 2 of 2 each (\S\ref{sec:frozen})\\
H200 cell, 300k pool & none & & & window $(1,\,1.3)\,\lref$ too narrow: the rule's rates 1.10 and 1.20 fall outside its margin band $(1.1,\,1.17)$, and the fill returns INCOMPLETE; no test\\
\fpath{2026-10-30_llama31.yaml} (Llama-3.1-8B E2) & FILLED & \texttt{1559f8e} 10-31 07:49:13 & 07:49:15 & 18 runs, submitted 07:50:02 (17170409--26); window $(1,\,2.0)\,\lref$, 1.35 and 1.65 $\lref$, cap 20; pass at 1.65 $\lref$, fail at 1.35 (\S\ref{sec:frozen})\\
\fpath{2026-10-31_h200_p600k.yaml} (H200 600k E2) & FILLED & \texttt{4d1d431} 10-31 09:26:38 & 09:26:41 & 18 runs, submitted 09:27:44 (17171020--37); window $(1,\,2.5)\,\lref$, 1.5 and 2.0 $\lref$, cap 48; pass at both rates (\S\ref{sec:frozen})\\
\bottomrule
\end{tabular}
\end{table*}

\noindent Superseded before registration: a native rescue engaged at 2400\,s, whose simulator prediction (13 of 20 rescued) led to the engage time of 2100\,s that R4 registers (\fpath{notes/PREDICTIONS_r11_rescue.md}). The SHA-256 of every tagged file is in its tag message (\texttt{r11-basin-reg}, \texttt{r11-lpminf-pred}, \texttt{r11-rescue-reg}, \texttt{r14-basin2-reg}, \texttt{r14-extra-reg}) and in \fpath{provenance/registration-proofs/}; the Llama-3.1-8B and 600k H200 registrations are not tagged, and the job script writes their SHA-256 into each of their runs' \fpath{meta.json}.

\section{Prediction rounds and timeline}\label{app:timeline}

\paragraph{The calibrated simulator in detail.} Our discrete-event simulator models a continuous-batching engine: iterations of at most 8192 prefill tokens and 256 sequences; admission reserves a request's end-of-turn KV; idle KV of sessions in a tool call or waiting is evicted whole, least recently used first; waiting requests pin nothing; a finished session's KV is freed at once, an oracle, since a real engine keeps it as evictable cache. An iteration costs a fixed time plus a per-token term plus attention terms. The five cost parameters were fitted only to SGLang's per-batch timing lines in the Stage~0 and Stage~C server logs (decode $R^2=0.992$, median error 0.75\%; prefill $R^2=0.998$, 1.3\%), never to $\mu(n)$. It runs mixed prefill--decode iterations or, like SGLang, separate iterations with prefill priority; the primary (mixed) mode was the one closer to the real Stage~C $\mu(n)$ under a pre-declared metric (mean relative error 0.104 against 0.127), a margin carried by one cell, and both modes' predictions are reported. It does not model partial prefixes, the waiting-prefix refresh or the overlap scheduler; reserving whole turns at admission is an upper bound on what paged engines hold (in round 5 of the earlier simulator, on-demand allocation with preemption and SGLang-like ratio reservation kept the curve and the trap).

\paragraph{Validation.} The declared level test, simulated $\mu$ within 15\% of the real value in each of 8 Stage~C cells, failed in both modes: mixed batching missed one cell, FCFS $n=48$, by $+42\%$ (0.070 against 0.049\sps), separate batching two. The shape was reproduced: the hit fraction crosses 0.5 at $n=32.5$ against 32.2, the plateau ($n=16$--24) agrees within 9.2\% and $n=300$ within 2\%, and $\mu(300)/\mu^*$ is 0.50 against 0.46. The single real run at $n=48$ dipped below the simulator's all-miss level; its turns-based estimate, 0.063\sps, is within 15\%. Keeping a finished session's KV as evictable cache, a labelled sensitivity arm, left mixed batching the closer mode (0.118 against 0.146) but lowered the plateau ($-15.5\%$ at $n=24$; 5 of 8 cells within 15\%). In the earlier simulator the same change moved the orderings and pins of \S\ref{sec:mech} by at most 0.025 $\lmiss$ but stalled the CONCUR surrogate at its minimum cap in 12 of 18 cells; the static caps were not rerun.

\paragraph{Prediction rounds.} The calibrated simulator then predicted, in both modes (mixed primary), B1: $\lmiss$ and the all-miss token rate; B2: the fraction of from-empty runs that stay stable at each rate; B3: the E2 outcome at each rate from 0.09 to 0.16\sps; B4: the closed curve with 600k- and 900k-token pools. Two later rounds repeated these predictions for the native pool (B1n--B4n) and for vLLM (B1v--B4v, with the simulator's token costs scaled to vLLM's measured all-miss rate, a fit, not a prediction), and three more, for the tests of \S\ref{sec:frozen}, were made by the same simulator (rounds 11 and 12) and by the first-principles model of \S\ref{sec:fp} (round 13, in two sets: tag \texttt{r13-theory-pred}, and \texttt{r13-theory-pred2} for BASIN-2, TOOLSLOW, SHED and the H200 600k cell, with the pool-sweep lookups and the rival scoring of the 0.5-hit crossing used for the pool sweep, the 300k H200 and Llama-3.1-8B). LLM coding agents made the simulator predictions, each in its own copy of the repository and instructed not to read the results they predicted; every round was committed and pushed under a tag before its results were read by its authors. What that means differs by round: B3, B4 and B3n are blind in time; B2 was pushed after E0b had finished on the cluster but before its results were pulled, and B2n, B2v and B4v after their jobs were submitted but before any had finished; B1, B1n and B4n are not blind, and B1 is close to the calibration data; B3v predates the vLLM registration, but its BURST part is blind only on the predicting agent's side. An independent verifier reran the first round. Appendix~\ref{app:timeline} gives the details and the time record of every round.

\begin{table}[p]
\centering\footnotesize
\caption{Order of events, 2026-10-21 to 10-26 (times CDT). Commits are those of the development repository (availability statement); commit and tag times are recorded by the committer, push times by GitHub's activity log, job times by Slurm.}\label{tab:timeline}
\begin{tabular}{@{}p{1.75cm} >{\raggedright\arraybackslash}p{5.6cm}@{}}
\toprule
When & Event\\
\midrule
10-21 & Stage 0 passes (17060165); Stage C submitted 15:28 (17062103--05)\\
10-22 12:54 & Stage C results pulled\\
10-22 13:13 & runbook proposes the rate rule and a cap of 20 (\texttt{288b801})\\
10-22 13:14 & E0a submitted (17072634--37); the last job ends 16:34\\
10-22 13:50 & rate rule and cap adopted in the status notes; simulator cost model fitted to Stage 0/C server logs (\texttt{4435c51})\\
10-22 17:03 & Stage 0 and C data committed (\texttt{d0235a7})\\
10-22 17:45 & first push to GitHub, with the three commits above\\
10-23 14:32 & E0a results committed (\texttt{a76f5e5})\\
10-23 14:42 & E0b submitted (17084750--55); jobs end by 16:09\\
10-24 03:13 & validation of the calibrated simulator committed (\texttt{813cfc1})\\
10-24 03:23 & prediction spec and code committed before any prediction run (\texttt{630ab9e}); on GitHub 03:24\\
10-24 03:49 & predictions B1--B4 committed (\texttt{64bc4e3})\\
10-24 04:05 & independent verification recorded (\texttt{877b936})\\
10-24 14:52 & tag \texttt{r6-predictions-frozen} (tagger time)\\
10-24 14:53 & B1--B4 and the verification reach GitHub\\
10-24 15:06 & E0b results pulled\\
10-24 15:09 & E0b data and registration committed in the same second (\texttt{cfbc702}, \texttt{d279b7b}); registration on GitHub 15:09:54\\
10-24 15:13 & E2 submitted (17093387--404)\\
10-24 15:38 to 10-25 02:36 & E2 jobs run\\
10-25 02:28 & native-pool Stage C and E0a start (17096856--58)\\
10-25 14:40 & registered E2 verdict computed (data \texttt{49fbc7b})\\
10-25 14:42 & native Stage C curve analysed on the cluster; native curve and $\lmiss$ in the status notes at 14:47 (\texttt{6793fb7})\\
10-25 14:45 & native E0b (17100938--39) and vLLM E0a (17100940) start\\
10-25 14:49 & native prediction spec committed (\texttt{29728e3})\\
10-25 14:52 & native Stage C and E0a data committed (\texttt{fa265c2})\\
10-25 14:54 & B1n--B4n run on the clean code commit \texttt{0336c91}\\
10-25 15:06 & B1n--B4n committed (\texttt{e3ea48f}); on GitHub 15:08:28\\
10-25 15:08 & tag \texttt{r8-native-predictions-frozen} (tagger time); native rate rule in the status notes (\texttt{6cb43e9})\\
10-25 17:26 & native E0b jobs end; pulled 17:30\\
10-25 17:31 & native E0b data and native registration committed (\texttt{f3f55ab}, \texttt{b8c19ec}); registration on GitHub 17:31:56\\
10-25 17:42 to 10-26 01:11 & native E2 jobs (17102173--90)\\
10-25 18:06 & vLLM calibration and exploratory registration committed (\texttt{65b65e2}, \texttt{a78ce07}), on GitHub 18:06:33; vLLM BURST submitted 18:07 (17102318--23)\\
10-25 22:42 to 10-26 03:24 & vLLM BURST jobs\\
10-26 10:07 & native E2 and vLLM BURST verdicts computed (data \texttt{f116cf3})\\
\bottomrule
\end{tabular}
\end{table}

\begin{table}[htbp]
\centering\footnotesize
\caption{Order of the vLLM events on 2026-10-30 (times CDT, recorded as in Table~\ref{tab:timeline}).}\label{tab:timeline-vllm}
\begin{tabular}{@{}p{1.75cm} >{\raggedright\arraybackslash}p{5.6cm}@{}}
\toprule
When & Event\\
\midrule
12:10 & vLLM Stage C (17162959--61) and E0b (17162962--63) submitted; 17162959--60 start 12:11\\
12:16 & rule for the vLLM cap and rates in the status notes, on GitHub (\texttt{836063a})\\
12:31 & B1v--B4v committed (\texttt{cc01dbf}); tag \texttt{r9-vllm-predictions-frozen} 12:32 (tagger time)\\
12:32 & B1v--B4v reach GitHub (12:32:44)\\
14:36 & first vLLM jobs end (Stage C, 17162959--60)\\
15:35 & vLLM Stage C curve analysed on the cluster\\
16:23 & bracket of the five finished vLLM E0b runs (16:23:19); vLLM registration committed (\texttt{5a6fbcd}, 16:23:59)\\
16:24 & paired test submitted (17166121--38, 16:24:36); first job starts 16:25:42, second 16:26:17\\
16:26 & vLLM registration and Stage C data (\texttt{276c5f0}) reach GitHub (16:26:35)\\
17:17 & vLLM E0b job 17162963, with the 3.0 $\lref$ run, ends\\
17:28 & vLLM E0b data committed and pushed (\texttt{582c2a6})\\
\bottomrule
\end{tabular}
\end{table}

\noindent Commit and tag times are set by the committer. The external record is GitHub's activity log for the development repository, archived in \fpath{provenance/github_activity_2026-10.json} and, refreshed after the last registered job, \fpath{provenance/github_activity_2026-11-01.json}, together with the cluster's accounting records of the cited jobs (\fpath{provenance/sacct_cited_jobs_2026-10-30.txt}, \fpath{provenance/sacct_cited_jobs_2026-11-01.txt}). The development repository is private; the archived log is an unsigned export made by the first author with GitHub's API, and the hash chains in \fpath{provenance/registration-proofs/} show that the files are unchanged, not when they were pushed; \fpath{provenance/push_times_2026-10.tsv} and its refresh \fpath{provenance/push_times_2026-11-01.tsv}, which adds the rounds 11 to 14, list, for each registration and prediction commit, its first arrival on GitHub and the first job it concerns. GitHub records no push time for tags, so the push times given here are those of the first branch update that contains each commit. The branch holding B1--B4 (\texttt{64bc4e3}) and the verification reached GitHub at 14:53:07 CDT on 10-24, and the E2 registration (\texttt{d279b7b}) at 15:09:54, before the E2 jobs were submitted at 15:13:39. On 10-25, B1n--B4n (\texttt{e3ea48f}) arrived at 15:08:28, the native registration (\texttt{b8c19ec}) at 17:31:56 and the exploratory vLLM registration (\texttt{a78ce07}) at 18:06:33. Pull times are from our local records only.

B3 (E2) and B4 (native pool) reached GitHub before their experiments were submitted, so they are blind in time. B2 (E0b) was committed after E0b had finished on the cluster but before its results were pulled. The LLM agents that made the predictions worked in separate copies of the repository at a commit without the E0a or E0b results and were instructed not to contact the cluster and not to open the E0a or E0b result directories or the primary copy of the repository; they could technically have done so. Their tool-call transcripts, which are not in the repository, were checked afterwards by the orchestrating Claude Code session, which found no cluster command and no access to E0b before the predictions reached GitHub. For B2 this instruction-level control, checked after the fact, is the only safeguard. A scan of every line added between the commit the agents started from and the predictions found no use of withheld data, but a text scan cannot prove that nobody read it. B1 (E0a) rests on the same instruction only: E0a was committed before the predictions. The independent verifier, a separate OpenAI Codex session, had read the project status notes, which mention E0a, so it is independent but not blind; it also reused the prediction agents' copy of the repository instead of a fresh one, as its record states. Its evidence is in \fpath{provenance/codex_prediction_verification/}. Two conditions of the verification failed and are disclosed: the strict containment of B3 passes within the B2 all-stable range, and the platform-specific test failures of the full suite.

In the second round, B2n (native E0b) and B3n (native E2) reached GitHub while native E0b was still running, before any of its runs had finished, and before the native registration or any native E2 job existed. B3n is therefore blind in time. B2n reached GitHub 23\,min after the two native E0b jobs were submitted, when both were running, so it is blind in the sense that no result existed, not in submission order. B1n and B4n are not: native Stage C and E0a had finished, and their headline numbers were in the project status notes before the round's specification was written. That specification is shared by all four, so B2n and B3n are blind to the native E0b and E2 outcomes but were specified after the native curve was known, and B3n's cap of 64 is B4n's output. The LLM agent that ran the round worked in a separate copy of the repository at a commit that predates the native data (\texttt{fb3510e}) and, by its own record, read nothing under the real-engine results directories; that record has not been checked independently. No independent verifier reran the second round.

In the third round (\fpath{notes/PREDICTIONS_vllm.md}), an LLM agent in a separate copy of the repository at \texttt{a4bab1d}, which holds the exploratory vLLM BURST data but no vLLM closed curve or edge, ran the calibrated simulator with its token-proportional costs scaled by one factor per batching mode (1.0882 mixed, 1.0806 separate), chosen so that its all-miss token rate equals vLLM's measured 10,349 tokens/s; since the factor is fitted, B1v is not a prediction. By its own record, which has not been checked independently, it read nothing under the real-engine results directories; the two real vLLM numbers it used, $\lmiss$ and the all-miss token rate, were given in its task. The predictions (\texttt{cc01dbf}, committed 12:31:58) reached GitHub at 12:32:44, 21\,min 48\,s after the five vLLM Stage~C and E0b jobs were submitted (Table~\ref{tab:timeline-vllm}); two of these jobs had been running since 12:11 and none ended before 14:36:50. So B2v and B4v are blind in the sense that no result existed when they were pushed, not in submission order; B3v predates the vLLM registration and the paired test. The orchestrating Claude Code session that specified the round knew the outcome of the exploratory vLLM BURST runs, which were in the repository by then, and chose the candidate rates, so the BURST-arm predictions at 1.35 and 1.65 $\lref$ are blind at most on the predicting agent's side. The notes' remark that their branch was never pushed predates the push. No independent verifier reran the third round. The vLLM registration's own timing is in Appendix~\ref{app:prereg}.

\paragraph{Rounds 11 to 14.} Table~\ref{tab:timeline-r11} gives the time record of the later rounds. Round 11 froze the mechanism checks for the registered vLLM test (\S\ref{sec:vllm}) and the simulator's predictions for the basin, fresh-seed, LPM-INF and rescue tests; round 12 the simulator's predictions for the Llama-3.1-8B cell, with its token costs scaled to the cell's measured all-miss rate (a fit); round 13 the first-principles model of \S\ref{sec:fp}; round 14 registered the BASIN-2, TOOLSLOW and SHED tests. LLM agents ran the simulator rounds in separate copies of the repository. The round-11 rescue predictions were revised after two adversarial reviews and before any rescue run, and every registered quantity was bit-identical before and after the revision (\fpath{notes/PREDICTIONS_r11_rescue.md}). The agents that derived the first-principles model read every committed result of the three measured cells and the registered protocols of BASIN-1, R1 and R4, and no H200, Llama, pool-sweep, basin, fresh-seed, LPM-INF or rescue result (\fpath{notes/PREDICTIONS_r13_theory.md}, ``Blindness scope''). When they froze the basin and rescue predictions they also knew the simulator's burst-size sweep of \S\ref{sec:robust} (bursts of 25 never trapped, 50 about half the time and 100 always) and the registered simulator rescue prediction (R1 rescued in 20 of 20 seeds at both rates). Their only H200 inputs were its E0a measurements. The agent that operated the H200 cell had shown its closed-curve values for $n\le32$ to the first author before the freeze; the agents that derived the model did not see them.

Table~\ref{tab:timeline-r11} also gives what each prediction is blind to. The simulator's BASIN-1 and fresh-seed predictions reached GitHub before those jobs were submitted, and its LPM-INF predictions before the LPM-INF jobs, so all three are blind in time. For the theory, no BASIN-1 run is blind in time. The three 25-session runs had started 1.7 to 2.2\,h before the push (6230--7807\,s of 8700\,s elapsed, past the burst and into the window [5700, 8700)\,s), and the three 100-session runs had ended 0.2 to 1.6\,h before it, so for all six the prediction is blind only by access: no one is known to have read them before the push, and the first recorded reading of any BASIN-1 or fresh-seed outcome was at about 04:07, but no file-access record can show this, because the cluster's class storage is mounted without access times. The same holds for the H200 cell's closed runs at $n=40$, 48, 96 and 300 and its original E0b runs. They had ended before the freeze, and they were first read after it, when the curve and the bracket were combined at 02:30. The runs at $n=36$, 44 and 52 were submitted after the freeze and are blind in time. No blind simulator prediction exists for the H200 cell. The second theory set (pred2) reached GitHub at 04:20:20, before the BASIN-2, TOOLSLOW and SHED jobs were submitted (04:21:40 to 04:21:43), so its predictions for them are blind in time. The 600k H200 jobs had been submitted at 02:50 but could not start before 05:00 and first started at 05:27, so its 600k predictions are blind in time by the start criterion of Table~\ref{tab:frozensum}: no result existed at the push. The pool-sweep jobs had been submitted at 22:58 on 10-30, before either set was pushed, but were held until 02:30 and first started at 06:50:13, so both sets are blind in time for them by the same criterion. The model's a priori Llama-3.1-8B predictions (pushed at 02:05:30) and the simulator's (02:14:06) are blind by access only for that cell's closed runs, which ran from 00:50 to 03:31: three had ended at the theory push, three were running and three started after it (from 02:26); by the operator's record none was read before both pushes (the curve was first combined from 03:14, for $n=16$--24, and in full by 03:37). That cell's from-empty runs started after 03:40 and are blind in time for both, as are its paired-test jobs (submitted at 07:50:02, before that push; first start 12:41) for the registration-conditional predictions pushed at 07:52:42 and 07:52:56. pred2's table for the fresh seeds is not blind: it was committed at 04:20:05, after the last fresh-seed run had ended (03:57:56) and after those outcomes were first read (about 04:07); by its record its agents read no result, which is why its own erratum calls it blind by access only. The simulator's rescue predictions (pushed at 21:30:40 on 10-30) and the model's (02:05:30 on 10-31) reached GitHub before the first rescue run started (02:23:44 on 10-31), so all nine rescue runs are blind in time to both; the smoke run ended before the model's push, but the model made no prediction for it, and by their record the agents that derived the model had no Nova access and opened no rescue output before the push.

\begin{table*}[tbp]
\centering\footnotesize
\caption{Order of the round-11 to round-14 events, 2026-10-30 and 10-31 (times CDT; commit and tag times by the committer; submission, start and end times of jobs by Slurm's accounting, run end times by each run's \fpath{meta.json}). Push times marked * were read from GitHub's activity log at the time (recorded in \fpath{notes/PREDICTIONS_r11_mechanism.md}, \fpath{notes/PREDICTIONS_r13_theory.md}, and \fpath{notes/PREDICTIONS_r13_pred2.md}) or on 10-31; each agrees with the export of 2026-11-01 (\fpath{provenance/github_activity_2026-11-01.json}; rows of \fpath{provenance/push_times_2026-11-01.tsv}, except \texttt{f1338e3} and \texttt{d12eef2}, which are in the export only). The last rows run past 02:00 CDT on 11-01, when daylight saving time ended; times after it are CST.}\label{tab:timeline-r11}
\begin{tabular}{@{}p{2.3cm} >{\raggedright\arraybackslash}p{14.4cm}@{}}
\toprule
When & Event\\
\midrule
10-30 18:20 & r11 mechanism predictions for the vLLM test committed (\texttt{f1338e3}, 18:20:11) and pushed (18:20:13*); three of its 18 jobs running, none ended\\
10-30 18:55 & first vLLM paired-test job ends (17166121, 18:55:07)\\
10-30 19:50 & BASIN-1 and fresh-seed predictions (\texttt{c600fd7}); BASIN-1 registrations (\texttt{9128995}, 19:51); tag \texttt{r11-basin-reg} 20:04; all pushed 20:04:30*\\
10-30 21:01 & LPM-INF predictions (\texttt{e2e8474}) and registration (\texttt{5064580}, 21:02); tag \texttt{r11-lpminf-pred} 21:02; pushed 21:02:46*\\
10-30 21:12 & BASIN-1 and fresh-seed jobs submitted (17168500--07, 21:12:45--46)\\
10-30 21:28 & rescue predictions (\texttt{72501f0}); tag \texttt{r11-rescue-pred} 21:30; pushed 21:30:40*\\
10-30 22:07 & LPM-INF jobs submitted (17168743--45, 22:07:40)\\
10-30 22:10 & R1 and R4 registrations (\texttt{f0bc3ad}); tag \texttt{r11-rescue-reg} 22:38; pushed 22:39:02*\\
10-30 22:28 & the H200 operator reads the closed-curve values for $n\le32$ (17166818--22) and shows them to the first author\\
10-30 22:38 to 23:52 & H200 closed runs at $n=40$, 48, 96 and 300 end (17166828--34; not read until 02:30)\\
10-30 22:58 & pool-sweep jobs submitted (17168864--81, 22:58:03); eligible to start from 02:30:32 on 10-31\\
10-30 23:01 & rescue jobs submitted (17168965--74, 23:01:12--13); the smoke job runs 23:24 to 00:08\\
10-31 00:14 to 01:08 & H200 E0b jobs end (17168377--82; not read until 02:30)\\
10-31 00:20 & last vLLM paired-test job ends\\
10-31 00:28 to 01:51 & the three 100-session BASIN-1 runs end (00:27:56, 00:29:27, 01:50:54)\\
10-31 00:50 to 03:31 & Llama-3.1-8B closed runs (17166814--16): three end before the theory push, three start after it (from 02:26); combined from 03:14 ($n=16$--24) to 03:37 (all nine)\\
10-31 01:58 & BASIN-2 registrations (\texttt{621d88f}); tag \texttt{r14-basin2-reg} 02:10; pushed 02:10:48*\\
10-31 02:05 & theory frozen (\texttt{997d2c4}, 02:04:55) and pushed (02:05:30*); tag \texttt{r13-theory-pred}\\
10-31 02:05 to 03:58 & the six fresh-seed runs end (the first at 02:04:59, 31\,s before the theory push)\\
10-31 02:08 & registered vLLM verdict computed on the cluster; pulled 02:09:56, committed 02:22 (\texttt{cacf833})\\
10-31 02:14 & Llama simulator predictions (\texttt{a62aea9}) pushed (02:14:06*); tag \texttt{r12-llama31-predictions-frozen}\\
10-31 02:20 to 02:47 & the three 25-session BASIN-1 runs end\\
10-31 02:23 to 06:57 & the nine registered rescue runs (17168966--74; the first starts at 02:23:44, 18 minutes after the theory push; the last ends at 06:57:06)\\
10-31 02:27 & TOOLSLOW and SHED registrations (\texttt{b006291}); tag \texttt{r14-extra-reg} 02:27:59; pushed 02:28:09*\\
10-31 02:30 & H200 nine-point curve and E0b bracket combined on the cluster: first reading of $n\ge40$ and of the edge\\
10-31 02:33 & H200 closed runs at $n=36$, 44 and 52 submitted (17169612--14, 02:33:06); they end 03:25 to 04:01\\
10-31 02:50 & H200 E0b reseeds (17169668--71, 02:50:20) and 600k jobs (17169672--80, 17169682--87, 02:50:32 to 02:51:20, not eligible to start before 05:00) submitted\\
10-31 03:49 to 04:21 & the three LPM-INF runs end\\
10-31 04:07 & first recorded reading of the BASIN-1 and fresh-seed outcomes (about 04:07)\\
10-31 04:20 & second theory set for BASIN-2, TOOLSLOW, SHED and the H200 600k cell, with the pool-sweep lookups and the rival scoring of the 0.5-hit crossing, committed (\texttt{cfa7712}, 04:20:05) and pushed (04:20:20*); tag \texttt{r13-theory-pred2}\\
10-31 04:21 & BASIN-2, TOOLSLOW and SHED jobs submitted (17169893--911, 04:21:40 to 04:21:43) from clones at the registration tags, by scripts that require the pred2 commit\\
10-31 05:27 & first 600k H200 job starts (17169672--80, 82--87); curve and bracket combined at 09:24\\
10-31 06:50 to 11:04 & pool-sweep runs (17168864--81); pulled at 11:23 and scored against both sets (\texttt{3025139})\\
10-31 07:05 & rescue runs pulled (committed in \texttt{c6ea41e}); registered verdicts computed at 07:17 (\texttt{be50593}, \texttt{454ff7a})\\
10-31 07:49 & Llama-3.1-8B registration (\texttt{1559f8e}, 07:49:13) pushed (07:49:15*); its 18 E2 jobs submitted at 07:50:02 (17170409--26)\\
10-31 07:52 & registration-conditional Llama predictions (\texttt{a36f441}, \texttt{d12eef2}) pushed (07:52:42*, 07:52:56*), before any Llama E2 job started (first start 12:41)\\
10-31 09:26 & 600k H200 registration (\texttt{4d1d431}, 09:26:38) pushed (09:26:41*); its 18 E2 jobs submitted at 09:27:44 (17171020--37)\\
10-31 09:27 to 11-01 00:28 & the 18 600k H200 paired-test runs (17171020--37)\\
10-31 10:10 to 15:11 & the nine BASIN-2 runs (17169893--901); pulled at 15:42 and first read then\\
10-31 12:41 to 23:31 & the 18 Llama-3.1-8B paired-test runs (17170409--26)\\
10-31 21:03 to 11-01 01:00 CST & the six TOOLSLOW and four SHED runs (17169902--11)\\
11-01 15:44 to 15:49 CST & registered analyses over all runs, with the frozen code: the 600k H200 and Llama-3.1-8B verdicts computed on the cluster and pulled (15:45); the TOOLSLOW and SHED runs pulled and analysed locally (\texttt{db24efa}, 15:49)\\
\bottomrule
\end{tabular}
\end{table*}

\subsection{Blind predictions against outcomes}\label{sec:blind}

\begin{table}[!tp]
\centering\footnotesize
\caption{Blind predictions (calibrated simulator, mixed batching unless noted) against outcomes; rates in sessions/s; native and vLLM rates in multiples of $\lmiss$ (for the simulator its own, for the engine that cell's $\lref$). Percentages in the last column are predicted/real $-1$. B1 is not blind in time and is close to the calibration data (\S\ref{sec:sim}, Appendix~\ref{app:timeline}). B3 and B3n used a per-arm pass rule (each arm in 2 of 3, resp.\ 7 of 10, seeds); the registered rule asks for all three arms in the same seed, which B3n's runs also meet (9 of 10 seeds at its low mixed point, 10 of 10 elsewhere). B3's token-ratio and backlog ranges are over the two simulated rates that bracket each registered rate. B4 is the registered seed~0 on the grid $n=48$, 64, 96 (supplementary 5-seed means: peak 0.514, $\mu(96)=0.085$, $\mu(300)=0.108$). B2n and B3n: 10 seeds (``sep.'': separate batching); B3n's low and high points are 2.85 and 4.65 (sep.\ 2.65 and 4.35). B4v and B2v: vLLM, variant V1 (simulator token costs scaled to vLLM's all-miss rate), 10 seeds, turns-based rates. B3n is blind in time; B2n, B4v and B2v were pushed after their jobs were submitted, before any run had finished (\S\ref{sec:vllm}); B4n (3-seed mean) is not blind. B3v: vLLM paired test, all four configurations (V0, V1; mixed, separate), 10 seeds; blind in time, its BURST part blind only on the predicting agent's side (\S\ref{sec:vllm}).}\label{tab:blind}
\setlength{\tabcolsep}{2.5pt}
\begin{tabular}{@{}l l l@{}}
\toprule
 & predicted & real\\
\midrule
B1 & $\lmiss$ 0.0722 (turns) & 0.0715 ($+0.9\%$)\\
   & 11,220 tokens/s & 11,243 ($-0.2\%$)\\
   & $\lmiss$ 0.0694 (sessions) & 0.0717 ($-3.1\%$)\\
\addlinespace[1pt]
B2 & all-stable edge 0.14 (3 seeds), 0.15 (20) & in $[0.143,\,0.179)$ (seed 0)\\
   & at 0.18: 2/3 stable (15/20) & 0.179: 3/4 seeds stable\\
\addlinespace[1pt]
B3 & PASS at all rates 0.09--0.16 & PASS at 0.0966, 0.1180\\
   & \quad 20 seeds: sep.\ fails at 0.16 & \quad (0.16 not run)\\
   & token ratio 1.003--1.027 (0.09--0.10) & 1.009--1.025 (0.0966)\\
   & \quad 1.025--1.043 (0.11--0.12) & 1.025--1.029 (0.1180)\\
   & backlog 171--234 (0.09--0.10) & 194--226 requests\\
   & \quad 253--313 requests (0.11--0.12) & 277--325\\
   & good-arm hit 0.929--0.937 & 0.921--0.934\\
   & trapped 3/3 at 0.09 and 0.10; & trapped 2/3 at 0.0966\\
   & \quad 20 seeds: 14/20, 18/20 & \\
\addlinespace[1pt]
B4 & peak $n=64$, $\mu$ 0.523 (900k) & $n=64$, 0.479 ($+9\%$)\\
   & 0.5-hit crossing $\approx$80 & $\approx$93 ($\approx$91 on B4's grid)\\
   & $\mu(96)=0.079$, $\mu(300)=0.109$ & 0.123, 0.101\\
\midrule
B2n & all-stable edge 6.5; sep.\ 6.0 & in $[5.5,\,6.5)$\\
    & at 6.5: 10/10 stable; sep.\ 8/10 & 6.5 unstable\\
    & at 1.5--5.5: 10/10 stable, both & stable\\
\addlinespace[1pt]
B3n & PASS at 2.85, 4.65; sep.\ 2.65, 4.35 & PASS at 2.5, 4.0\\
    & trapped token ratio 1.086--1.109 (low) & 1.074--1.078 (2.5)\\
    & \quad 1.119--1.140 (high) & 1.106--1.124 (4.0)\\
    & trapped backlog 460--643 (low) & 549--594 requests\\
    & \quad 1067--1333 requests (high) & 1089--1130\\
    & good-arm hit 0.928--0.936 & 0.929--0.935\\
    & BURST trapped 9/10 to 10/10 & 3/3 at both\\
\addlinespace[1pt]
B4n & peak $n=72$, $\mu$ 0.532 & 64--80, 0.479 ($+11\%$)\\
    & 0.5-hit crossing 94; hit 0.42 at 96 & $\approx$93; 0.42\\
    & cap 64, 0.35\% inside the rule & 64\\
\midrule
B4v & peak $n=24$, $\mu$ 0.215; cap 20 & $n=20$, 0.211 ($+2\%$); 20\\
    & hit 0.74 at $n=28$, 0.46 at 32 & 0.66, 0.08\\
    & 0.5-hit crossing 31.4 & 29.1\\
\addlinespace[1pt]
B2v & all-stable edge 2.0; at 2.5: 7/10 & in $[2.0,\,2.5)$; 2.5 unstable\\
\addlinespace[1pt]
B3v & PASS at 1.35, 1.65 (all four) & PASS at both\\
    & BURST trapped 7--9/10 & 2/3 at each\\
    & NOBURST, CAPPED(20) 10/10 & 3/3 and 3/3 at each\\
    & trapped token ratio 0.99--1.07 & 0.948--1.018\\
\bottomrule
\end{tabular}
\end{table}

B1 matched within 3.1\% (Table~\ref{tab:blind}), which mostly checks the simulator's bookkeeping (\S\ref{sec:sim}). B2's edge, 0.14--0.15\sps, is consistent with the nominal seed-0 bracket $[0.143,\,0.179)$; near 0.179 the model keeps 2 of 3 runs stable (15 of 20 with more seeds), so one unstable run there is no contradiction, and the engine was stable there on the three fresh seeds of \S\ref{sec:edge} (3 of 4 seeds stable in all). B3 predicted a pass at every tested rate in both batching modes. At the simulated rates that bracket each registered rate, the primary (mixed) mode's trapped token ratios, 1.003--1.027 at 0.09--0.10\sps{} and 1.025--1.043 at 0.11--0.12, contain the observed 1.009--1.025 at 0.0966 and 1.025--1.029 at 0.1180; but at 0.12, the grid point nearest the higher rate, the prediction was 1.034--1.043, above what the engine did. The observed good-arm hit fractions (0.921--0.934) are close to, and at the low end partly below, the predicted 0.929--0.937. B3 flagged both ends of its grid as fragile. At the low end the slope margin $0.5(\lambda-\lmiss)$ is small: with its registered three seeds it trapped 3 of 3 at both 0.09 and 0.10, but with 20 seeds 14 and 18, so 10--30\% of runs there are stuck without meeting the slope test; one indeterminate real run in three at 0.0966 is consistent with that. At 0.16 the 20-seed supplement of the secondary (separate) mode failed the pass rule, with NOBURST good in 13 of 20 seeds, one short; the engine was not run there. (Backlogs compare only roughly: the simulator counts waiting requests, the engine runs outstanding ones.) B4 got the peak's position, within its coarse grid, and the outward shift of the cliff; the predicted peak was 9\% above the real one and the cliff early: on B4's grid of 64 and 96 the real hit fraction crosses 0.5 near 91, against B4's 80, and at $n=96$ the engine still hit 42\% of turns, against B4's 5\%. The early cliff is mostly the pool: B4n, the same simulator at the real pool (not blind), hit 0.42 at $n=96$ and crossed 0.5 near 94. The peak error is not: at $n=64$, where nothing is evicted, seed 0 gives the same 0.523 in both pools. Two qualifications: the verifier's strict check that every B3 pass lies inside the B2 all-stable range failed (B3 asks for 2 of 3 seeds over a longer run), and B2 is blind only by instruction (checked afterwards from the agents' transcripts) and by commit order, not by event order: E0b had finished on the cluster before the predictions were committed, and was pulled 13 minutes after their commit reached GitHub.

The second round, for the native pool, held in part: B3n's pass held at the registered rates, which lie 6--16\% below its own test points, with trapped token rates, backlogs and good-arm hit fractions consistent with the lower rates; the primary mode of B2n put the from-empty edge at least one grid step too high; and B1n and B4n are not blind (Appendix~\ref{app:supp}). For vLLM, B4v and B2v got the peak level, the cap and the edge right but predicted a later and gentler cliff, and B3v's pass held (\S\ref{sec:vllm}).

\paragraph{Slips in the registration record.} The vLLM registration reached GitHub after the first two of its jobs had started (\S\ref{sec:prereg}), and the memory request of its 15 pending jobs was lowered after submission (Appendix~\ref{app:e2detail}). B1n and B4n were specified after the native curve was known, and B2n, B2v and B4v were pushed after their jobs were submitted (Appendix~\ref{app:timeline}). The pool-sweep jobs (for both theory sets), the 600k H200 jobs (for pred2) and the Llama-3.1-8B paired-test jobs (submitted at 07:50:02, for the registration-conditional predictions pushed at 07:52:42 and 07:52:56) were also submitted before the predictions were pushed, but none started until after the push, so we count those predictions as blind in time by the start criterion of Table~\ref{tab:frozensum}. The agent that operated the H200 cell saw its closed-curve values for $n\le32$, which include the predicted peak, before the freeze and showed them to the first author; the agents that derived and froze the model did not see them, so the H200 peak-ratio rule and hit$(32)$ are not blind to the first author, while the runs at $n=36$, 44 and 52 are blind in time and the other runs at $n\ge40$ and the original from-empty runs, which had ended before the freeze, are blind by access only; four H200 E0b reseeds were added after the first H200 E0b had been read (\S\ref{sec:frozen}). No blind simulator prediction was made for the H200 cell. The 100-session runs of BASIN-1 had ended 0.2--1.6\,h before the model's predictions were pushed, and the 25-session runs were then 6230--7807\,s into their 8700\,s; no one is known to have read them before then, but the cluster's file system records no access times, so all six predictions are blind by access only. The H200 runs at $n=36$, 44 and 52 and the fresh-seed from-empty runs ended before pred2 was pushed and were read before it (combined at 04:09; first read at about 04:07), so pred2's tables for them are not blind, nor are its rows for BASIN-1's burst sizes, whose outcomes the agent that operated the basin cell also first read at about 04:07; pred2's agents read no BASIN-1 result, so to them those rows are blind by access only. The freeze record of pred2, \fpath{notes/PREDICTIONS_r13_pred2.md}, says it was pushed before the first job of the H200 runs at $n=36$, 44 and 52 (those jobs were submitted at 02:33 and had ended before the push); its dated erratum (\texttt{bf19385}) corrects this and calls both tables blind by access only, counting only its agents' access. The BASIN-2, TOOLSLOW and SHED jobs were submitted after the push (04:20:20 CDT; jobs submitted 04:21:40--43). pred2's rival scoring of the 0.5-hit crossing, including the test of the tool-call term, was declared after the 300k H200 and Llama-3.1-8B curves had been combined by their operators (the H200 curve at 02:30; the Llama curve from 03:14, in full by 03:37), so for those two cells it is not blind. The model's and the simulator's Llama-3.1-8B predictions were pushed while that cell's closed runs were in progress, so for its curve they are blind by access only; and the Llama and 600k H200 registrations were pushed 47 and 63\,s before their jobs were submitted and are not tagged, nor are the Llama registration-conditional predictions (\texttt{a36f441}, \texttt{d12eef2}). The agents that froze the basin and rescue predictions of the model knew the simulator's burst-size sweep and its registered rescue prediction. The engage time of the native rescue test was chosen after a simulator prediction for an earlier engage time had been seen, before any rescue run. The bands of the frozen vLLM mechanism checks were chosen on SGLang runs and on exploratory vLLM runs whose arrivals the registered BURST arm repeats, so only its no-burst and capped arms are out of sample for those checks. No registration or prediction is in a third-party registry, and none has a timestamp other than GitHub's push log of the development repository, which is private, so readers see that log only in our export.

\section{Provenance}\label{app:prov}
Every number in the paper is listed with its source file and commit or job in \fpath{paper/NUMBERS.md}. Table~\ref{tab:prov} gives the data directory and commit behind each result. The commit hashes cited are those of the development repository; the public repository's history is a redacted copy of it, and \fpath{provenance/registration-proofs/} maps each cited registration and prediction commit to its public counterpart (availability statement). The 103 commits made before 17:45 CDT on 2026-10-22, including \texttt{288b801}, \texttt{4435c51} and \texttt{d0235a7} cited here, predate a history rewrite on that day that removed the AI co-author trailers from their commit messages and changed nothing else: trees, authors, dates and parent relations are unchanged. \fpath{provenance/commit-map.tsv} maps their old hashes, which notes and run records written before the rewrite cite, to the current ones; the repository the old hashes belong to is not on GitHub, so they can be checked only against this map. Their earliest server record is the repository's first push to GitHub, at 17:45:50 CDT on 2026-10-22, with which 97 of them arrived, these three among them, before E0b was submitted on 2026-10-23; one more (\texttt{fc870c2}) arrived 2\,s later, and five are reachable only from archive tags, which have no push record.

\paragraph{Repository, redactions and timing evidence.} Its history is a redacted copy of the authors' development repository: passages unrelated to the paper were removed from four notes files and one clause from each of two commit messages, five local tool-configuration files were dropped from every commit (all listed in \fpath{provenance/REDACTIONS.md}), and no code, experiment configuration, registration, result or paper file changed in it. A final release commit adds licences and README text and replaces allocation names in some notes, job-script comments and tests (\fpath{provenance/}). Its commit hashes therefore differ from those cited in this paper, which are the development repository's. The directory \fpath{provenance/} maps every cited hash to its public counterpart and holds the timing evidence: GitHub's push log of the development repository (\fpath{github_activity_2026-10.json}, refreshed in \fpath{github_activity_2026-11-01.json}), the external record of when each registration and prediction reached GitHub (\S\ref{sec:prereg}, Appendices~\ref{app:prereg} and~\ref{app:timeline}); for every registration and frozen prediction, the original commit, tree and tag objects, the push-log event that first carried them, and a script that recomputes the hash chains and checks that each registration file in the public repository is byte-identical to the registered one (\fpath{provenance/registration-proofs/}); for the registrations and predictions pushed after 2026-10-30, the refreshed push log (\fpath{github_activity_2026-11-01.json}) and the first push of each with the first job it concerns (\fpath{push_times_2026-11-01.tsv}); the cluster's accounting records of all 258 cited jobs (\fpath{sacct_cited_jobs_2026-11-01.txt}, and in UTC \fpath{sacct_cited_jobs_2026-11-01_utc.txt}, exported after the last job had ended; the earlier \fpath{sacct_cited_jobs_2026-10-30.txt} covers the jobs through the vLLM paired test); and the evidence of the independent verifier of the first prediction round. Simulator outputs are not stored, apart from their summaries; each can be regenerated from the command, commit and configuration recorded with the run (\fpath{meta.json}, \fpath{config.yaml}). The derived tables in \fpath{paper/derived/} name each engine run by its path on the cluster, \fpath{results/agentload_*/<timestamp>}; the run is stored under the pull directory of its jobs in \fpath{results/nova/}, with the same final path components. The agents' tool-call transcripts (\S\ref{sec:sim}, Appendix~\ref{app:timeline}) are not in the repository.

\begin{table*}[tbp]
\centering\footnotesize
\caption{Data behind each result (paths relative to the repository root at tag \texttt{arxiv-v1}).}\label{tab:prov}
\begin{tabular}{@{}>{\raggedright\arraybackslash}p{2.0cm} >{\raggedright\arraybackslash}p{14.7cm}@{}}
\toprule
Result & Directory (commit)\\
\midrule
Stage 0, C & \fpath{results/nova/s0_smoke_17060165_pulled_20261021-152347/}, \fpath{results/nova/stageC_1706210x_pulled_20261022-125454/} (\texttt{d0235a7})\\
E0a & \fpath{results/nova/e0a_1707263x_pulled_20261023-142037/} (\texttt{a76f5e5})\\
E0b & \fpath{results/nova/e0b_1708475x_pulled_20261024-150639/} (\texttt{cfbc702})\\
E2 & \fpath{results/nova/e2_170933xx_pulled_20261025-144558/} (\texttt{49fbc7b})\\
native, grid, vLLM & \fpath{results/nova/r7_native_grid_vllm_pulled_20261025-145126/} (\texttt{fa265c2})\\
native E0b & \fpath{results/nova/e0b_native_1710093x_pulled_20261025-173038/} (\texttt{f3f55ab})\\
vLLM E0a & \fpath{results/nova/e0a_vllm_17100940_pulled_20261025-180814/} (\texttt{b273241})\\
native E2, vLLM BURST & \fpath{results/nova/e2_native_vllm_burst_pulled_20261026-101256/} (\texttt{f116cf3})\\
vLLM C & \fpath{results/nova/stageC_vllm_1716295x_pulled_20261030-153621/} (\texttt{276c5f0})\\
vLLM E0b & \fpath{results/nova/e0b_vllm_1716296x_pulled_20261030-172810/} (\texttt{582c2a6})\\
validation & \fpath{results/r6_validate_stagec/20261024-025222/} (\texttt{813cfc1})\\
B1--B4 & \fpath{results/r6_predict_report/20261024-034707/}, \fpath{notes/PREDICTIONS.md} (\texttt{64bc4e3})\\
B1n--B4n & \fpath{results/r8_native_report/20261025-150146/}, \fpath{notes/PREDICTIONS_native.md} (\texttt{e3ea48f})\\
B1v--B4v & \fpath{results/r9_vllm_report/20261030-122723/}, \fpath{notes/PREDICTIONS_vllm.md} (\texttt{cc01dbf})\\
vLLM E2 & \fpath{results/nova/e2_vllm_1716612x_pulled_20261031-020956/} (\texttt{cacf833}); mechanism checks in its \fpath{r11_mechanism/} and \fpath{derived/}\\
theory & \fpath{results/theory/merged/}, \fpath{notes/PREDICTIONS_r13_theory.md}, \fpath{notes/THEORY_r13_scoring.md}, \fpath{configs/r13_theory_predict.yaml} (\texttt{997d2c4}, tag \texttt{r13-theory-pred}); \fpath{results/theory/paper/}, \fpath{figures/theory/} (\texttt{5271ff6})\\
H200 & \fpath{results/nova/h200_wave1_pulled_20261031-023416/} and \fpath{prereg/inputs/2026-10-31-h200/} (\texttt{d108056}); \fpath{results/nova/h200_stageC_extra_1716961x_pulled_20261031-040933/} (\texttt{748900b}); \fpath{results/nova/h200_e0b_reseeds_1716966x_pulled_20261031-050308/} (\texttt{4e27364}); 600k pool \fpath{results/nova/h200p600k_stageC_e0b_171696xx_pulled_20261031-092822/} (\texttt{8b548c0}), registration \fpath{configs/agentload_prereg_2026-10-31_h200_p600k.yaml} (\texttt{4d1d431})\\
BASIN-1, E0F & \fpath{results/nova/b1_e0f_1716850x_pulled_20261031-040517/} (\texttt{4dd21e0})\\
BASIN-2 & \fpath{results/nova/b2_1716989x_pulled_20261031-154500/}; scoring \fpath{experiments/r14_b2_score.py}, \fpath{experiments/r14_b2_derived.py} (\texttt{eb08a5d})\\
pool sweep & \fpath{results/nova/pool_sweep_1716886x_pulled_20261031-112342/}, scored by \fpath{experiments/r13_pool_score.py} into its \fpath{derived/} (\texttt{3025139})\\
Llama-3.1-8B & \fpath{results/nova/llama31_wave1_pulled_20261031-033750/} (\texttt{c4c9d50}); \fpath{results/nova/llama31_e0b_registration_pulled_20261031-075313/} (\texttt{38f0712}); conditional predictions \fpath{results/theory/pred2/conditional/} (\texttt{a36f441}, \texttt{d12eef2})\\
LPM-INF & \fpath{results/nova/lpminf_1716874x_pulled_20261031-042548/}, \fpath{results/r11_lpminf_eval/20261031-eval/} (\texttt{fc8b028})\\
rescue R1, R4 & \fpath{results/nova/rescue_1716896x_pulled_20261031-070502/} (\texttt{c6ea41e}); verdicts \fpath{results/agentload_analysis_rescue_r1/20261031-071720/} (\texttt{be50593}) and \fpath{results/agentload_analysis_rescue_r4/20261031-071748/} (\texttt{454ff7a}) in it; descriptive tables \fpath{results/rescue_paper_data/20261031-071938/} in it (\texttt{51c803d}); registrations \fpath{configs/agentload_prereg_2026-10-31_rescue*.yaml} (\texttt{f0bc3ad}); predictions \fpath{results/r11_predict_rescue/20261030-195231/} (\texttt{72501f0})\\
Proposition~\ref{prop:E} & \fpath{results/theory/pred2/prop_e_cor_c.json} (\texttt{cfa7712}, tag \texttt{r13-theory-pred2})\\
pred2 & \fpath{results/theory/pred2/}, \fpath{configs/r13_theory_pred2.yaml}, \fpath{notes/PREDICTIONS_r13_pred2.md} (\texttt{cfa7712}, tag \texttt{r13-theory-pred2})\\
rounds 11, 12 & \fpath{results/r11_predict_basin/20261030-194539/} (\texttt{c600fd7}), \fpath{results/r11_predict_lpminf/20261030-194911/} (\texttt{e2e8474}), \fpath{results/r11_predict_rescue/20261030-195231/} (\texttt{72501f0}), \fpath{results/r12_llama31_report/20261031-021057/} (\texttt{a62aea9})\\
gap$/R_M$ & \fpath{paper/derived/loopvars_*.csv} (frozen script of \texttt{f1338e3})\\
robustness & \fpath{results/r7_report/20261024-233554/} (\texttt{c653498}); \fpath{results/r7_traces_sim/r7_report/20261025-012451/}, \fpath{results/r7_traces_sim_adj/r7_report/20261025-012453/} (\texttt{b530e55})\\
mechanisms & \fpath{results/prefix_meta_baselines_audit_summary/20261025-001850/} (\texttt{dc76032}); \fpath{results/prefix_meta_r4_curve/20261021-050856/}, \fpath{results/prefix_meta_r4/20261021-051550/} (\texttt{a1d485b}); \fpath{results/prefix_meta_mech/20261021-025902/} (\texttt{18ff187})\\
push log, jobs & \fpath{provenance/}\\
figures & \fpath{figures/paper/}, \fpath{figures/rescue_paper_data.py}, \fpath{paper/figs/README.md}\\
derived & \fpath{paper/derived/}, \fpath{figures/paper/derived_numbers.py}\\
\bottomrule
\end{tabular}
\end{table*}

\paragraph{Derived numbers.} Quantities that no committed analysis reports directly are computed from committed run data by \fpath{figures/paper/derived_numbers.py} and written to \fpath{paper/derived/}: the CAPPED drain times; the BURST backlog after the burst and over the window; the E0b hit fractions and realized arrival rates; the arm counts and the post-hoc sign test that counts each seed once; per BURST run, the hit fraction before the burst and the first 100\,s window from which no turn hit; and, per run of every arm, the client-side cost (gate waits, turn latency, unfinished sessions) and the realized arrival and completion rates over the window. They cover E0b, E2, native E0b, native E2, vLLM E0b and the exploratory vLLM BURST runs. Each derived file names a run both by its path on the cluster and by its path from the repository root.

\paragraph{Notes that support the text.} The math check of the model, \fpath{notes/MATHCHECK_thrashing_curve.md}, was done by an LLM agent on 2026-10-21; it predates the present assumptions (A1)--(A4) and Lemma~\ref{lem:waits}, and where it differs from \S\ref{sec:model} and Appendix~\ref{app:proofs}, those hold; the verification of Stage~C (three-source hit check, eviction survival times, prompt continuity) is \fpath{notes/STAGEC_VERIFICATION.md}; the audit of the PEEK, Continuum and CONCUR rules against released code is \fpath{notes/BASELINES_AUDIT.md}; an audit of the citations as of 2026-10-25, the scoop check (Lu et al.~\citep{lu2026cachetrap} appeared on 2026-11-02, after the last registered jobs of the earlier version had run (about 2026-11-01), and this version positions our work against it; \fpath{notes/LIT_CHECK_trap_2026-11-05.md} holds the literature check for model~F) and a first reading of the SGLang and vLLM sources are in \fpath{notes/lit/CITATION_AUDIT_2026-10-25.md} (that reading went twice through a page summariser, a web tool that returns a language model's summary of each file, without verbatim code; the line references of \S\ref{sec:background} and Appendix~\ref{app:proofs} were later read verbatim in the files at the tags), and that audit predates the citations added since, among them Yuan et al.~\citep{yuan2026agentic}, which it lists as not cited; the trace statistics are in \fpath{notes/lit/AGENT_TRACES.md}. The theory's derivations and their review are \fpath{notes/THEORY_r13_merged.md}, \fpath{notes/THEORY_r13_fluid_rigorous.md} and \fpath{notes/THEORY_r13_che_endogenous.md}; where they differ from \S\ref{sec:fp} and Appendix~\ref{app:fp}, those hold. Where the research log and a committed CSV differ in the last digit (three hit fractions and one backlog of the policy grid), the paper uses the CSV.

\section{Proofs of the quasi-static results}\label{app:proofs}

\paragraph{The congestion bound.} Section~\ref{sec:props} states the following proposition in brief.

\begin{proposition}[Request-level FCFS, any online eviction rule]\label{prop:A}
Consider the closed system with $n$ sessions under any online eviction rule, one that does not foresee tool times, and assume that its long-run rates exist and are deterministic (Appendix~\ref{app:proofs}) and that:
\begin{enumerate}[label=(A\arabic*),leftmargin=*]
\item sessions share no cached tokens, each has at most one outstanding request, and evicted KV is discarded, not offloaded to another memory tier; tool times are i.i.d.\ with finite mean, independent of all other randomness, and unknown to the engine until they end;
\item admission is request-level FCFS: no request is admitted for the first time before a request that was already waiting when it joined the queue. More generally, it suffices that the order meets the condition of Appendix~\ref{app:proofs} for some $Q_0$ and $S_0$ independent of $n$, which allows reordering while at most $Q_0$ requests wait and lets a follow-up pass at most $S_0$ requests that were waiting when it joined. At most $B$ requests run at once, and at most $\alpha t+\beta$ admissions, re-admissions after a preemption included, occur in any interval of length $t$, with $B$, $\alpha$ and $\beta$ independent of $n$;
\item in congestion the engine computes prompt tokens no faster than with caching disabled: $\limsup_{n\to\infty}u_n\le C$, where $u_n$ is the long-run rate of prompt tokens that are computed rather than served from cache when their request is first admitted (each prompt token is counted once, so tokens recomputed after a preemption are not), and $C$ is the limit of that rate as $n\to\infty$ for the closed system with caching disabled.
\end{enumerate}
Then the long-run rates of cache-served prompt tokens and of hit turns tend to 0 as $n\to\infty$, so does the hit fraction if $\liminf_{n\to\infty}\mu(n)>0$, and $\limsup_{n\to\infty}\mu(n)\le\lmiss$. If also (A4) $u_n\to C$, that is, the congested engine stays saturated and computes prompt tokens as fast as with caching disabled, then $\mu(n)\to\lmiss$, so the hit fraction tends to 0.
\end{proposition}
\noindent The proof below rests on Lemma~\ref{lem:waits}. The idea is that a follow-up can reuse a cached token only if the token stays in the pool while the follow-up waits, which bounds reuse by $M$ over the wait whatever the eviction rule, and that request-level FCFS makes short waits rare as $n$ grows, because the engine cannot foresee tool times.

\begin{remark}[resident-first orders; a conjecture]
Under pure LPM (without SGLang's 128-request fallback) in congestion we expect the follow-up hit fraction to tend to $P(Z<R_\infty)$, with $R_\infty$ the congested lifetime of idle KV, so $\muinf^{\mathrm{LPM}}$ exceeds $\lmiss$ through tool calls shorter than $R_\infty$; the simulator agrees (\S\ref{sec:mech}), but we have no proof. That service in arrival order sustains a metastable state is known from other systems~\citep{bronson2021metastable,farahbakhsh2026characterizing}, and serving the newest requests first once a queue forms is deployed practice~\citep{maurer2015fail}; favouring a job whose cache footprint is still resident is classical affinity scheduling~\citep{squillante1993affinity}; FCFS and LPM under prefix reuse are compared analytically by Dexter et al.~\citep{dexter2025prefix}; and PEEK documents the 128-request fallback~\citep{xie2026peek}. The exploratory LPM-INF test of \S\ref{sec:frozen} disables that fallback on the engine.
\end{remark}


\paragraph{Model and conventions.} In the closed system of \S\ref{sec:notation}, $n$ sessions are present and a finished session is replaced at once by a fresh one, whose first request joins the queue. At time 0 every session is fresh, so no tool call is in progress, as in the synchronized start of our closed runs. A request joins the queue when its session sends it, waits, and is admitted; it leaves the queue only by admission. An admitted request runs until it completes or, if the engine preempts it, rejoins the queue. At a common instant joins are processed before admissions, as when an engine receives requests before each scheduling pass, and requests that join at the same instant enter in some order; the requests waiting when request $i$ joins are those in the queue just before it joins and those that join at the same instant ahead of it. When a request that is not its session's last completes, the session starts a tool call, and when the call ends the session's next request, a follow-up, joins the queue. By (A1) evicted KV is discarded, not offloaded. A follow-up's wait runs from its joining the queue to its first admission, and its reuse, and whether it hits, are counted at that admission. SGLang logged no retraction in any of the 73 server logs of our SGLang runs. When a running request cannot get a block for its next tokens, vLLM preempts the most recently admitted running request and returns it to the head of its queue (\emph{Other orders} below); its closed runs logged 12 to 72 preemptions per run from $n=32$ on (72 at $n=300$, in a run that completed 5,395 requests), and its cache-disabled reference 9 or 10 per 32-session run. Since $u_n$ and $C$ count each prompt token once, at its request's first admission, recomputation after a preemption enters neither, and the proofs below go through with re-admissions counted among the admissions that (A2) bounds. By (A1) each session is at every moment in exactly one of three states: its request waits, its request runs, or it is in a tool call. Tool times have law $F_Z$ and mean $\bar z<\infty$. A long-run rate is $\lim_{t\to\infty}A(t)/t$ for a count or sum $A(t)$ over $[0,t]$; we assume that these limits exist almost surely and are deterministic, as for a positive Harris recurrent Markov process from any initial state, and that $\lmiss$ and $C$, the limits of $\mu(n)$ and $u_n$ as $n\to\infty$ with caching disabled, exist. Every $A(t)/t$ used below is bounded for fixed $n$ and $t\ge1$, since at most $\alpha t+\beta$ requests are admitted by time $t$, so by bounded convergence each rate also equals $\lim_{t\to\infty}E[A(t)]/t$, and bounds on expectations carry over to rates.

\begin{lemma}[short waits are rare]\label{lem:waits}
Assume (A1) and (A2), with $Q_0$ and $S_0$ as in (A2) (both 0 for request-level FCFS). Let $\gamma_n(K)$ be the long-run rate of follow-ups that wait less than $K$, and let $\kappa=\lfloor\alpha+\beta\rfloor+B$, $y=\kappa(1+\bar z)$ and $m=n-B-\alpha K-\beta-Q_0-S_0$. If $m>\kappa$, then
\[
\gamma_n(K)\le\frac{y+y^2+\kappa y}{m-\kappa},
\]
so $\gamma_n(K)=O(1/n)$ for every fixed $K>0$.
\end{lemma}

\noindent The general ordering condition of (A2), the only property of the order that the proof uses, is this: if more than $Q_0$ requests wait at every moment from a follow-up $i$'s joining the queue to its first admission, counting at a moment with admissions the requests that wait just before them, then all but at most $S_0$ of the requests that were waiting when $i$ joined are admitted, for the first time or again, no later than that admission. Request-level FCFS meets it with $Q_0=S_0=0$, since $i$ itself waits.

\begin{proof}
For a follow-up $i$ write $q_i$ for the time it joins the queue, $D_i$ for its wait and $Q_i$ for the number of requests waiting when it joins, and let $X(t)$ be the number of sessions in tool calls at time $t$.

\emph{A short wait needs a short queue.} Let $D_i<K$. At most $\alpha K+\beta$ admissions occur in $[q_i,q_i+D_i]$, $i$'s own included, so fewer than $\alpha K+\beta$ occur before $i$'s. If $Q_i\ge\alpha K+\beta+Q_0+S_0$, then at every moment of $[q_i,q_i+D_i]$, counted as above, more than $Q_0$ of the $Q_i$ requests still wait, because requests leave the queue only by admission; so at least $Q_i-S_0$ of them are admitted no later than $i$, and $Q_i-S_0+1\le\alpha K+\beta$, a contradiction. Hence $Q_i<\alpha K+\beta+Q_0+S_0$.

\emph{A short queue needs almost every session in a tool call.} Just before $q_i$ each of the other $n-1$ sessions is waiting (at most $Q_i$ of them), running or in a tool call, and $i$'s own session is in a tool call or, if that call has length 0, running; at most $B$ sessions run. So $X(q_i-)\ge n-Q_i-B>m$ whenever $D_i<K$.

\emph{From the instants at which tool calls end to fixed times.} Follow-ups join the queue when tool calls end, which happens more often when many sessions are in tool calls, so an average over these instants is not a time average; we avoid it as follows. A request that completes in an interval $(t_0,t_0+1]$ was admitted, for the last time, in that interval or was running at time $t_0$, so at most $\kappa$ requests complete, and at most $\kappa$ tool calls start, in any such interval. Fix an integer $j\ge0$. A follow-up that joins the queue in $(j,j+1]$ ends a tool call that was in progress at $j$ or started in the interval, so at most $X(j)+\kappa$ follow-ups join in it, and $X(q-)\le X(j)+\kappa$ for every $q\in(j,j+1]$. By the previous step, the number of follow-ups that join in $(j,j+1]$ and wait less than $K$ is at most $g(X(j))$, where $g(x)=(x+\kappa)\,\mathbf 1\{x>m-\kappa\}$.

\emph{A bound on $X(t)$ that does not depend on $n$.} Fix $t$ and split $(-\infty,t]$ into the intervals $I_k=(t-k-1,\,t-k]$, $k\ge0$; no tool call starts before time 0, and at most $\kappa$ start in each $I_k$. Let $(Z_{kl})_{k\ge0,\,1\le l\le\kappa}$ be i.i.d.\ with law $F_Z$, independent of everything else. Give the $l$-th call that starts in $I_k$ the duration $Z_{kl}$, and calls that start after $t$ fresh durations. This does not change the law of the system: by (A1) each call's duration is independent of all other randomness, and the system's evolution up to any time depends on a call's duration only through whether the call has ended by then; the index $(k,l)$ of a call is fixed when the call starts, so the variable it receives has not been used and is independent of the past at that moment, and the two constructions agree in law by induction over successive events. A call that started in $I_k$ and is still in progress at $t$ has a duration greater than $k$, so in this construction
\[
X(t)\le\hat X=\sum_{k\ge0}\sum_{l=1}^{\kappa}\mathbf 1\{Z_{kl}>k\},
\]
whose law depends on neither $t$ nor $n$. $\hat X$ is a sum of independent indicators with $E[\hat X]=\kappa\sum_{k\ge0}P(Z>k)\le\kappa(1+\bar z)=y$, so $\mathrm{Var}\,\hat X\le E[\hat X]$ and $E[\hat X^2]\le y+y^2$.

\emph{Conclusion.} Since $g$ is non-decreasing and $\mathbf 1\{\hat X>m-\kappa\}\le\hat X/(m-\kappa)$,
\[
E[g(X(j))]\le E[g(\hat X)]\le\frac{E[\hat X^2]+\kappa E[\hat X]}{m-\kappa}\le\frac{y+y^2+\kappa y}{m-\kappa}.
\]
Summing over the unit intervals of $(0,t]$, for integer $t$, bounds the expected number of follow-ups that join by time $t$ and wait less than $K$ by $t$ times this quantity; dividing by $t$ and letting $t\to\infty$ gives the claim.
\end{proof}

\noindent The bound is of order $1/n$, but its constants grow with the admission rate and the mean tool time, so it is a statement about the limit, not about the populations we ran.

\paragraph{Proof idea.} By (A1), the $r_i$ prompt tokens that follow-up $i$ reuses at its first admission are its session's own and stay in the pool from the end of the session's previous turn until that admission, an interval no shorter than its queue wait $D_i$, and the token-time that these intervals claim does not overlap. Resident token-time accrues at rate at most $M$, so follow-ups that wait at least $K$ reuse tokens at rate at most $M/K$, whatever the eviction rule. By (A2) a follow-up that waits less than $K$ joined behind fewer than $\alpha K+\beta+Q_0+S_0$ requests, so just before it joined more than $n-B-\alpha K-\beta-Q_0-S_0$ sessions were in tool calls; Lemma~\ref{lem:waits} (Appendix~\ref{app:proofs}), which uses that the engine does not see tool times, shows that such follow-ups occur at a rate of order $1/n$, and each reuses at most $L_{\max}=40{,}350$ tokens. Letting $n\to\infty$ and then $K\to\infty$, the cache-served token rate tends to 0, and since a hit reuses more than $L_0/2$ tokens, so does the rate of hits. Counted once, at its request's first admission, each of a session's $W$ prompt tokens is either served from cache or computed, so $u_n$ is $\mu(n)E[W]$ less the cache-served token rate, and with caching disabled the same count gives $C=\lmiss E[W]$; (A3) gives the bound and (A4) the limit. The same count gives turns at the rate $\mu(n)E[T]$, so the hit fraction tends to 0 whenever $\mu(n)$ stays bounded away from 0.

\begin{proof}[Proof of Proposition~\ref{prop:A}]
\emph{Reuse needs waiting token-time.} Let follow-up $i$ of session $\sigma$ join the queue at $q_i$, wait $D_i$, and reuse $r_i$ prompt tokens from cache at its first admission. By (A1) these tokens are $\sigma$'s own, computed by its earlier requests, and evicted KV is discarded, so each of them was in the pool throughout $[q_i,q_i+D_i]$, during which $\sigma$ has no running request to recompute it. These pairs of token and time are distinct for different requests: sessions share no tokens, and the waits of one session's requests do not overlap. At most $M$ tokens are resident at any time, so for every $t$, $\sum_i r_iD_i\le Mt$ over the follow-ups first admitted by $t$. This step uses neither the order nor the eviction rule, so it holds for clairvoyant rules too.

\emph{Reuse vanishes.} Fix $K>0$. By time $t$ the follow-ups that waited at least $K$ reused at most $Mt/K$ tokens, and each of the others at most $L_{\max}=40{,}350$, the longest prompt. So the long-run rate $\rho_n$ of cache-served prompt tokens satisfies $\rho_n\le M/K+L_{\max}\gamma_n(K)$, and Lemma~\ref{lem:waits}, with the order's $Q_0$ and $S_0$, gives $\limsup_{n\to\infty}\rho_n\le M/K$ for every $K$, so $\rho_n\to0$. A hit reuses more than half its prompt, hence more than $L_0/2$ tokens, so hits occur at a rate of at most $2\rho_n/L_0\to0$.

\emph{Work.} A session with $T$ turns sends $W(T)$ prompt tokens (\S\ref{sec:notation}), and each, counted once at its request's first admission, is either served from cache or computed. The closed system has $n$ slots, each holding one session at a time and a fresh one after it finishes, so each slot holds a sequence of sessions with i.i.d.\ turn counts, and by the strong law of large numbers completed sessions bring prompt tokens at the long-run rate $\mu(n)E[W]$ and turns at the rate $\mu(n)E[T]$; sessions in progress hold at most $nW(60)$ tokens and $60n$ turns and do not affect rates. Hence $u_n=\mu(n)E[W]-\rho_n$. With caching disabled nothing is reused, so the same count gives $C=\lmiss E[W]$. So $\mu(n)=(u_n+\rho_n)/E[W]$, and (A3) gives $\limsup_{n\to\infty}\mu(n)\le C/E[W]=\lmiss$. The hit fraction is the rate of hits over the rate of turns, so it is at most $2\rho_n/(L_0\,\mu(n)E[T])$ and tends to 0 whenever $\liminf_{n\to\infty}\mu(n)>0$. Under (A4), $\mu(n)\to\lmiss>0$.
\end{proof}

\paragraph{Scope of Proposition~\ref{prop:A}.} The proof's last step, from vanishing reuse to $\lmiss$, is the bound that MISA-T notes at prefill saturation, a request rate of at most the prefill throughput over the mean uncached prompt tokens per request~\citep{hong2026misat}; what the proposition adds is that request-level FCFS makes reuse vanish in congestion whatever the online eviction rule. (A2) describes an admission pass that stops at the first request that does not fit, as in our simulators and, on our reading of their sources, in SGLang and vLLM under their default FCFS policies (Appendix~\ref{app:proofs}); vLLM also preempts running requests and returns them to the head of its queue, which (A2) allows, because its ordering clause concerns first admissions; with asynchronous scheduling it also lets a follow-up pass a bounded number of just-preempted requests, so it meets the general condition of (A2) with $Q_0=0$ and a bounded $S_0$ (Appendix~\ref{app:proofs}). A diverging mean cycle time would not be enough: Lemma~\ref{lem:waits} uses bounded admissions and i.i.d.\ tool times to bound the tail of the number of sessions in tool calls independently of $n$, which makes it rare that almost all sessions are in tool calls when one returns. The lemma's bound is of order $1/n$ but has large constants, so the proposition describes the limit, not the populations we ran. The lemma needs FCFS order only while more than $Q_0$ requests wait, so, on our reading of its source, the proposition also covers SGLang's stock LPM, which stops re-sorting its queue whenever more than 128 requests wait, so that new requests queue behind all waiting ones ($Q_0=128$). It does not cover program-level FCFS, which puts an old session's follow-up ahead of newer sessions' requests however long the queue: in our earlier simulator its hit fraction stays near 0.03 from $n=64$ to 300, while FCFS has none. A pure resident-first order, such as LPM without the fallback, violates (A2) outright: a returning session whose KV is still resident is admitted ahead of the queue.

(A1) holds for our random-token sessions, apart from rare one-token coincidences (5 of the 19,924 first turns in E2 had one cached token), but not for agents that share a system prompt: one resident copy then serves every session. We expect the argument to apply to the private suffixes, with $\lmiss$ and $C$ measured with the shared prefix cached, but have not shown this. (A1) also rules out offloading evicted KV to host memory, which would let a follow-up reuse tokens that held no pool space while it waited. (A3) and (A4) hold on our engines only approximately: with zero hits, closed runs and the registered BURST runs computed prompt tokens at 0.95--1.12 times the cache-disabled rate $C$, and at 1.03 times it near the stationary turn mix, so under (A4) the proposition puts $\muinf$ about 2--5\% above $\lref$, within the standard error of $\lref$ (Appendix~\ref{app:proofs} gives the measurements).

The bound holds for every online eviction rule, including protection of queued prefixes and TTL pins that keep request-level FCFS order, where (A4) can fail; in our earlier simulator both still hit 10--23\% of turns at $n=300$, so a finite-$n$ hit fraction does not by itself separate covered orders from uncovered ones (Appendix~\ref{app:proofs}). Belady's offline rule~\citep{belady1966study} foresees tool times and is outside (A1): the token-time step of the proof covers it, Lemma~\ref{lem:waits} does not.

\paragraph{Checks of (A3) and (A4) on our engines.} (A3) and (A4) treat the engine as a server of prompt tokens, as E0a suggests (\S\ref{sec:edge}). They concern long-run rates of the closed system, in which every order serves turns in the stationary mix, since each completed session contributes all its turns; our measurement windows did not have that mix, and the token rate varied between them, so we can check (A3) and (A4) only approximately. We estimate $C$ by E0a's uncached token rate at 32 sessions, 11,243 tokens/s. With zero hits, closed runs computed prompt tokens at 1.03~$C$ at $n=96$, whose served prompts were within 3\% of the stationary mean, on SGLang and on vLLM (against vLLM's own $C$), at 0.99~$C$ at $n=48$ (prompts 12--14\% longer) and at 1.10--1.11~$C$ at $n=300$ (28--30\% shorter); the registered SGLang BURST runs computed 1.01--1.03~$C$ at 300k (prompts 9--20\% shorter) and 1.07--1.12 times the 32-session all-miss rate of the native pool (prompts 29--41\% shorter). So on our engine (A3) and (A4) hold only approximately: to within a few percent at the stationary mix, and to within 12\% over the windows we measured. We estimate $\lmiss$ by $\lref$, E0a's turns-based rate at 32 sessions (0.0715\sps); $\lref$ and $C/E[W]$ estimate the same cache-disabled limit. $C/E[W]$ is 0.0728\sps, 1.7\% above $\lref$, because the turns in E0a's window were longer than the stationary mix (10,648 against 10,469 prompt tokens on average). E0a's 64-session check computed 1.7\% more tokens per second than the 32-session runs and 3.2\% more than the same seed at 32 sessions, although its window held longer turns (11,456 prompt tokens on average), so the 32-session reference falls short of saturation by about 2--3\%, perhaps more. With the 64-session rate in place of $C$, $C/E[W]$ is 1.034 $\lref$, and with the zero-hit rate at $n=96$ (11,602 tokens/s) it is 1.050 $\lref$. Under (A4) the proposition therefore puts $\muinf$ about 2--5\% above $\lref$, an offset within the standard error of $\lref$ over seeds (0.0038\sps, 5\%); on vLLM the three estimates are 1.020, 1.038 and 1.050 $\lref$ ($\lref=0.0657$\sps).

\paragraph{Covered orders at finite $n$.} In our earlier simulator, under FCFS order, a hard pin of queued prefixes gives 0.919 $\lmiss$ at $n=300$ (\S\ref{sec:mech}) but 1.04--1.15 $\lmiss$ at $n=48$--150 (3 seeds), as a limit statement allows, and TTL pins of 5--20\,s, which that simulator keeps past expiry while the session's next request waits, give 0.10--0.29 $\lmiss$ at $n=300$ with the earlier release rule of \S\ref{sec:mech}, with the GPU busy only 8--33\% of the time (3 seeds). Both still hit 13\% and 10--23\% of turns at $n=300$, so a finite-$n$ hit fraction does not by itself separate covered orders from uncovered ones. The proposition makes their hit rate vanish as $n$ grows, and their hit fraction too unless their throughput tends to 0, as the TTL pins' may; the hard pin keeps the GPU busy, and its hit fraction falls from 0.23 at $n=48$ to 0.13 at $n=300$. The TTL rows of \S\ref{sec:mech} that exceed $\lmiss$ also serve pinned programs first, an order that (A2) excludes.

\paragraph{Other orders.} The proof uses the order only through Lemma~\ref{lem:waits}, so it holds for any order that meets the condition above for some $Q_0$ and $S_0$. In SGLang 0.5.22 the admission loop walks the waiting queue in order and stops at the first request that cannot be added or when the batch is full (\fpath{python/sglang/srt/managers/scheduler.py}, lines 4080--4176 at tag v0.5.22); it skips a request only for LoRA, pooling-mode or storage-prefetch conditions that our runs do not use. New requests join the queue at its tail (line 3366), as would a retracted request (lines 4369--4416). FCFS sorts the queue only under priority scheduling, which none of our runs enabled (\fpath{python/sglang/srt/managers/schedule_policy.py}, lines 227--231), and LPM sorts it only while at most 128 requests wait (lines 133 and 300--311); with more waiting, LPM takes the FCFS branch, which leaves the queue's order unchanged (lines 260--261), so new requests queue behind all waiting ones. So FCFS meets the condition with $Q_0=S_0=0$ and stock LPM with $Q_0=128$ and $S_0=0$. vLLM 0.32.0 admits from the head of its FCFS queue and stops at the first request whose allocation fails (\fpath{vllm/v1/core/sched/scheduler.py}, lines 1275--1282 at tag v0.32.0), and new requests join at the tail (\fpath{vllm/v1/core/sched/request_queue.py}, lines 75--94). It also preempts the most recently admitted running request when a running request cannot get a block (\fpath{vllm/v1/core/sched/scheduler.py}, lines 794--844), frees its blocks and returns it to the head of the queue (lines 1571--1613). The condition concerns first admissions, so this re-admission ahead of earlier requests does not violate it, and the re-admissions count among the admissions that (A2) bounds. With asynchronous scheduling, vLLM also skips a preempted request until its output still in flight drains, within the pipeline depth (lines 933--939), and a follow-up admitted meanwhile passes it; only requests preempted in the last few steps are skipped this way, at most a fixed multiple of the running limit, so $S_0$ stays bounded. vLLM also skips requests that are blocked, for example waiting for a KV transfer or over a LoRA limit, which our runs do not involve. Its FCFS thus meets the condition with $Q_0=0$ and a bounded $S_0$. The line numbers refer to the files at these tags, read verbatim, not through the page summariser of Appendix~\ref{app:prov}. Program-level FCFS meets the condition for no $Q_0$ and $S_0$: a follow-up of an old session is admitted before requests of newer sessions that were waiting when it joined, however long the queue.

\paragraph{Engine sources of \S\ref{sec:background}.} The SGLang rules of \S\ref{sec:background} (default policy and eviction rule, the 128-request LPM fallback, the unlocked prefix of a waiting request, and the refresh of waiting prefixes) are in these files under \fpath{python/sglang/srt/} at tag v0.5.22: \fpath{arg_groups/fields/schedule.py}, lines 82--98; \fpath{arg_groups/fields/memory.py}, lines 25--42; \fpath{managers/schedule_policy.py}, lines 133--135, 278--311 and 1591; \fpath{mem_cache/radix_cache.py}, lines 552--580; and \fpath{environ.py}, line 594. vLLM 0.32.0 schedules FCFS, its default, or by priority (\fpath{vllm/config/scheduler.py}, lines 21 and 159), and a freed cached block keeps its hash until it is reallocated (\fpath{vllm/v1/core/block_pool.py}, lines 666--690 and 774--803; at tag v0.32.0). We first read these sources through a web tool that returns a language model's summary of each file (Appendix~\ref{app:prov}); the line references here and above are to the files at the tags, read verbatim.

\paragraph{The quasi-static model in congestion.} Eq.~\eqref{eq:fluid} assumes that the internal state (turn-index mix, cache contents, queue order) relaxes quickly compared with changes in $n$. An automated check of the model by an LLM agent (Appendix~\ref{app:prov}) found this reasonable at fixed points and for the sign of the drift in deep congestion, where the stationary $\mu$ is flat, but not near the cliff. The model is also not quantitative in congestion, and not only transiently: over the evaluation window the registered SGLang BURST runs still completed 1.13--1.35 $\lref$ at 300k and 1.44--2.03 times the native $\lref$, and their fitted backlog slopes were only 18--69\% and 68--78\% of the rate $\lambda-\lmiss$ that Eq.~\eqref{eq:fluid} gives with $\mu$ at its congested value $\lmiss$, taking $\lambda$ as the arrival rate each run received and $\lmiss$ as $\lref$, or 21--75\% and 70--80\% with $\lmiss$ 3.4\% above $\lref$ (\S\ref{sec:e2}, \S\ref{sec:e2native} and the checks of (A3) and (A4) above). Under FCFS each waiting session advances one turn per cycle of the queue while new sessions keep arriving, so the backlog stays young and its turns are short (in the BURST windows served prompts were 9--20\% and 29--41\% below the stationary mean; checks of (A3) and (A4) above). At a token rate near the all-miss rate the engine then serves more turns per second than in the stationary mix, and since a session ends after each turn with probability 1/15 until its 60th, it also completes sessions faster than $\lmiss$, for as long as the backlog grows. Theorems~\ref{thm:trap} and~\ref{thm:premium} (\S\ref{sec:trap}) make this exact within model~F: the backlog grows at $\alpha=\lambda(1-E[z^{-T}])$ and sessions complete at $D=\lambda E[z^{-T}]>\lmiss$, with the turn mix tilted to $G_kz^{-k}$, and the premium does not fade while the trap lasts.

A work balance, which this paragraph does not make rigorous for the open system (Corollary~\ref{cor:trapthr} does, within model~F), suggests that the recovery threshold depends on the token rate in congestion, not on the departure rate. If hits vanish in congestion (as in every registered SGLang BURST run) and the engine computes prompt tokens at about $C=\lmiss E[W]$ (Proposition~\ref{prop:A}) while they arrive at $\lambda E[W]$, then for $\lambda>\lmiss$ the uncomputed prompt work of the sessions in the system grows, and eventually so does their number, since each holds at most $W(60)$ tokens of it, while for $\lambda<\lmiss$ that work shrinks as long as the engine computes at $C$. If the zero-hit token rate is lower at some population that the system passes through on the way down, the threshold is correspondingly lower; in the fluid model it is $\inf_{n\ge n^*}\mu(n)$ rather than $\muinf$. On our engine the closed runs' zero-hit token rate was 4\% lower at $n=48$ than at $n=96$, with older turns (checks of (A3) and (A4) above). We therefore use Eq.~\eqref{eq:fluid} for the window, as a long-run statement for the recovery threshold, and for the capped system's drain time (\S\ref{sec:props}), not for burst sizes or backlog growth in congestion: in an earlier simulator with assumed costs (\S\ref{sec:sim}) the fluid critical burst $n_2-n_1$ at $2\lmiss$ was 20.5 sessions against an empirical 40, because burst sessions are fresh and small. The heuristic H3 of \S\ref{sec:fp} corrects for this by counting a burst session at its peak footprint: at 1.65 $\lref$ on the 300k pool its critical fresh burst is 53.4 sessions, against 21.3 for the stationary-composition estimate $n_2-n_1$, and the registered basin tests of \S\ref{sec:frozen} test it.

\paragraph{A non-monotone tail, and Bronson et al.'s terms.} Without the monotone tail, the recovery threshold from deep congestion is $\inf_{n\ge n^*}\mu(n)$ instead, and congested fixed points can exist. Our measured curve is not monotone beyond $n^*$: in completed sessions it falls to 0.68 $\lref$ at $n=48$ and rises to 1.51 $\lref$ at $n=300$ (turns-based, 0.88 and 1.56 $\lref$; \S\ref{sec:curve}). Taken at face value, either estimate would put the recovery threshold from $n=300$ below $\lmiss$ and give a stable congested fixed point at 1.35 $\lref$. We attribute both features to the synchronized start, which no longer run has tested, and take $\muinf$ from Proposition~\ref{prop:A}. In the terms of Bronson et al.~\citep{bronson2021metastable}, $\muinf$ is the hidden capacity, the load up to which the system heals itself after any trigger, and $\min(\mu^*,\ledge)$ the advertised capacity; at loads in $\mathcal W$ the good state is their vulnerable state, a burst is a trigger whose intensity matters, the congested regime is the metastable failure state, and the feedback of \S\ref{sec:background}, in which misses lengthen the queue and the gap between a session's turns, is its sustaining effect.

\paragraph{The fluid statements of \S\ref{sec:fluid}.} With $\mu$ continuous and piecewise linear, Eq.~\eqref{eq:fluid} has a unique solution from every $x_0\ge0$; it is monotone in time, and it converges to a fixed point or diverges. Let $0<\lambda<\mu^*$. Then $\mu(n)=\lambda$ has exactly one solution $n_1$ in $(0,n^*)$, because $\mu$ increases strictly from 0 to $\mu^*$ there, and, since $\lambda\ne\muinf$, at most one in $(n^*,\infty)$, because there $\mu$ is non-increasing and strictly decreasing wherever $\mu>\muinf$. If $\lambda<\muinf$, then $\mu\ge\muinf>\lambda$ beyond $n^*$, so $\mu<\lambda$ on $[0,n_1)$ and $\mu>\lambda$ on $(n_1,\infty)$, and every solution converges to $n_1$. If $\muinf<\lambda$, the second solution $n_2$ exists, because $\mu(n^*)>\lambda$ and $\mu\to\muinf<\lambda$; then $\mu>\lambda$ on $(n_1,n_2)$ and $\mu<\lambda$ on $(n_2,\infty)$, so $n_2$ is unstable, and for $x_0>n_2$, $\dot n\ge\lambda-\mu(x_0)>0$ along the solution and $n\to\infty$. If $\lambda>\mu^*$, then $\mu<\lambda$ everywhere and every solution diverges. A solution from $x_0>n_1$ converges to $n_1$ if and only if $\mu>\lambda$ on all of $(n_1,x_0]$: if so, it decreases and stays in $(n_1,x_0]$, where $n_1$ is the only possible limit; if $\mu(x)\le\lambda$ for some $x\in(n_1,x_0]$, it never falls below $x$, where the drift is non-negative. For $x_0>n^*$, $\mu>\lambda$ on $(n_1,n^*]$ and the minimum of $\mu$ on $[n^*,x_0]$ is $\mu(x_0)$, so the condition is $\lambda<\mu(x_0)$, and $\mu(x_0)$ does not increase and tends to $\muinf$ as $x_0\to\infty$. Without the monotone tail the same argument gives the condition $\lambda<\min_{[n^*,x_0]}\mu$, which tends to $\inf_{n\ge n^*}\mu(n)$, and a point beyond $n^*$ where $\mu$ crosses $\lambda$ from below is a stable congested fixed point. Under the assumptions of Proposition~\ref{prop:A} this infimum is at most $\liminf_{n\to\infty}\mu(n)\le\lmiss$, and it equals $\lmiss$ if, in addition, $\mu(n)\ge\lmiss$ for every $n\ge n^*$.

\paragraph{Proof of Corollary~\ref{cor:rescue}.} While $N\ge c$ the admitted sessions form a closed system of $c$ sessions (Proposition~\ref{prop:cap}), so $\dot N=\lambda-\mu(c)$ and the gate, $N-c$, shrinks if and only if $\mu(c)>\lambda$. By unimodality $\mu>\lambda$ exactly on $(n_1,n_2)$, so the gate empties, from any initial population $N_0>c$, if and only if $n_1<c<n_2$; if $c\le n_1$ or $c\ge n_2$, then $\mu(c)\le\lambda$ and it never empties. Once $N<c$ the cap does not bind, $\dot N=\lambda-\mu(N)$ with $N<c<n_2$, and by the statements above $N$ stays in $[0,n_2)$ and converges to $n_1<c$, so the cap never binds again and removing it changes nothing. This is the standard closed-system equivalence of a system with a bounded multiprogramming level (\S\ref{sec:background}). If the closed curve itself folds (Proposition~\ref{prop:fold}), $\mu$ is multivalued on a band of populations and we take $n_2$ on the good branch; a cap $c\ge n_b$ engaged inside the trap keeps the all-miss state, which exists at every $n\ge n_b$, so $n_1<c<\min(n_2,n_b)$ is necessary, and we have not shown that it is sufficient. The stochastic counterpart (positive recurrence for $\lambda<\mu(c)$) should follow from standard arguments; we have not written them out.
